\documentclass[Afour,sagev,times]{sagej}

\usepackage[T1]{fontenc}
\usepackage[utf8]{inputenc}
\usepackage[spanish,es-tabla]{babel}
\usepackage{moreverb,url}
\usepackage{booktabs}
\usepackage{tabularx}
\usepackage{longtable}
\usepackage{array}
\usepackage{graphicx}
\usepackage{listings}
\usepackage{xcolor}
\usepackage{enumitem}
\usepackage{tikz}
\usetikzlibrary{shapes.geometric,arrows.meta,positioning,fit,backgrounds,calc}
\usepackage[colorlinks,bookmarksopen,bookmarksnumbered,citecolor=blue!50!black,urlcolor=blue!50!black,linkcolor=blue!50!black]{hyperref}

% Macros — kept to a minimum per Sage golden rule (i)
\newcommand{\code}[1]{\texttt{#1}}
\newcommand{\cq}[1]{\textsf{#1}}

% Chowlk-style TikZ helpers
\tikzset{
  daimoclass/.style={rectangle, draw=black!70, fill=yellow!20, thick, rounded corners=2pt,
    minimum width=2.4cm, minimum height=0.7cm, align=center, font=\ttfamily\scriptsize,
    inner sep=3pt},
  externclass/.style={rectangle, draw=black!60, fill=blue!10, thick, rounded corners=2pt,
    minimum width=2.4cm, minimum height=0.7cm, align=center, font=\ttfamily\scriptsize,
    inner sep=3pt},
  modulebox/.style={rectangle, draw=black, fill=gray!15, thick,
    minimum width=3.0cm, minimum height=0.9cm, align=center, font=\sffamily\footnotesize\bfseries,
    inner sep=5pt},
  vocabbox/.style={rectangle, draw=black!60, fill=green!12, thick, rounded corners=3pt,
    minimum width=2.2cm, minimum height=0.65cm, align=center, font=\sffamily\scriptsize,
    inner sep=3pt},
  subclassarrow/.style={-{Triangle[open,length=2.5mm]}, draw=blue!70!black, thick},
  objpropertyarrow/.style={-{Latex[length=1.8mm]}, draw=black!80, thin},
  modulearrow/.style={-{Latex[length=1.8mm]}, draw=black!60, thick, dashed},
}

\newcolumntype{L}[1]{>{\raggedright\arraybackslash}p{#1}}

\lstset{
  basicstyle=\ttfamily\footnotesize,
  breaklines=true,
  columns=fullflexible,
  showstringspaces=false,
  frame=single,
  framerule=0pt,
  backgroundcolor=\color{gray!10},
  captionpos=b,
  aboveskip=0.8em,
  belowskip=0.8em
}

\def\volumeyear{2026}

\begin{document}

\runninghead{Mori Orrillo y Liu}

\title{Un perfil de integración para el intercambio gobernado de modelos de IA en espacios de datos: la ontología DAIMO}

\author{Edmundo de Elvira Mori Orrillo\affilnum{1} y Jiayun Liu\affilnum{1}}

\affiliation{\affilnum{1}Departamento de Inteligencia Artificial, Escuela Técnica Superior de Ingenieros Informáticos, Universidad Politécnica de Madrid, España}

\corrauth{Edmundo de Elvira Mori Orrillo, Universidad Politécnica de Madrid, Campus de Montegancedo, 28660 Boadilla del Monte, Madrid, España.}

\email{edmundo.mori.orrillo@upm.es}

\begin{abstract}
El despliegue interorganizacional de modelos de inteligencia artificial exige un tratamiento semántico coherente de cinco tareas hoy fragmentadas entre vocabularios distintos: publicar el modelo como activo de catálogo, seleccionarlo bajo una política de uso, invocarlo mediante un servicio tipado, trazar cada ejecución y compararlo sobre condiciones de evaluación reproducibles. MLDCAT-AP describe el modelo, DCAT lo publica, ODRL expresa la política, PROV-O registra la procedencia y el \emph{Dataspace Protocol} (DSP) negocia el intercambio contractual; ninguno articula las cinco tareas como cadena única.

Este artículo presenta DAIMO (\emph{Dataspace AI Model Ontology}), un perfil de integración OWL~2~DL que reutiliza los vocabularios citados mediante alineamientos axiomáticos verificados. DAIMO introduce catorce clases nativas (nueve de primer nivel más cinco subclases de rol) para cubrir relaciones puente no expresadas directamente por los vocabularios reutilizados, y se acompaña de una capa SHACL con restricciones estructurales, seis invariantes SHACL-SPARQL entre clases y un banco de pruebas negativas.

DAIMO se desarrolla siguiendo la metodología \emph{Linked Open Terms} (LOT) y se ejercita sobre un demostrador multi-participante que instancia las veintitrés preguntas de competencia formuladas durante la especificación de requisitos. La evaluación reporta consistencia OWL, materialización OWL-RL, conformidad SHACL, pruebas negativas, consultas SPARQL y un benchmark sintético de escalabilidad acotada.
\end{abstract}

\keywords{gobernanza de modelos de IA; espacios de datos; alineamiento ontológico; validación SHACL; metodología LOT}

\maketitle

% =====================================================================
\section{Introducción}\label{sec:intro}

El Reglamento~(UE)~2024/1689, conocido como \emph{AI Act}, exige para los sistemas de inteligencia artificial de alto riesgo documentación técnica, registro de ejecuciones, trazabilidad y evidencia de auditoría~\cite{aiact}. Paralelamente, los espacios de datos europeos promueven intercambios contractuales entre participantes autónomos mediante catálogos, políticas y protocolos como el \emph{Dataspace Protocol} (DSP)~\cite{dsp}, implementado en conectores como Eclipse Dataspace Components (EDC)~\cite{edc}. En este contexto, un modelo de IA no es sólo un artefacto técnico: es un activo gobernado que debe poder publicarse, descubrirse, invocarse, trazarse y evaluarse a través de fronteras organizativas.

Los vocabularios existentes cubren fragmentos importantes de ese ciclo. MLDCAT-AP~3.0.0~\cite{mldcatap} extiende DCAT para describir modelos de aprendizaje automático como activos de catálogo, con metadatos sobre tareas, ejecuciones, evaluaciones, riesgos y recursos computacionales. DCAT~\cite{dcat2} aporta la capa de catálogo y servicios; ODRL~\cite{odrl22}, la representación de ofertas, permisos, prohibiciones y acuerdos; PROV-O~\cite{provo}, la procedencia de entidades, actividades y agentes; y DSP, el intercambio contractual en espacios de datos. Las \emph{model cards}~\cite{mitchell19}, \emph{factsheets}~\cite{arnold19} y \emph{datasheets}~\cite{gebru21} complementan estos vocabularios con documentación narrativa.

La brecha identificada en la literatura revisada es relacional. Un modelo descrito con MLDCAT-AP no queda, por ello solo, conectado al acto de ser ofrecido en un catálogo federado. Una política ODRL puede expresar permisos sin quedar ligada al despliegue que gobierna. Una negociación DSP no explicita necesariamente el contexto de evaluación ni la evidencia de auditoría asociada a una invocación de IA. Un bundle PROV puede agrupar procedencia sin distinguir la cadena que cruza varios contextos administrativos. El problema no es añadir otra descripción del modelo, sino articular catálogo, política, despliegue, ejecución, evidencia y evaluación como una cadena consultable y validable.

Este artículo presenta DAIMO (\emph{Dataspace AI Model Ontology}), un perfil OWL~2~DL de integración para modelos de IA gobernados en espacios de datos. DAIMO sigue una estrategia de reutilización preferente: conserva \code{it6:MachineLearningModel} como clase de modelo de MLDCAT-AP y utiliza DCAT, ODRL, PROV-O, DSP/EDC, FOAF, SKOS y SPDX allí donde ya proporcionan la semántica necesaria. Sólo introduce clases puente cuando la relación entre vocabularios no queda cubierta por esos recursos.

Las contribuciones se organizan en tres partes:

\begin{enumerate}[leftmargin=*]
\item un perfil de integración \emph{reuse-first} con catorce clases nativas, de las cuales nueve son clases de primer nivel y cinco son subclases de rol, orientadas a cubrir relaciones puente entre catálogo, política, despliegue, ejecución, procedencia y evaluación;
\item una evaluación reproducible del artefacto ontológico mediante consistencia OWL, materialización OWL-RL, verificación de implicaciones, conformidad SHACL, banco negativo de invariantes, escaneo OOPS!, veintitrés consultas de competencia y un benchmark sintético de escalabilidad acotada;
\item un demostrador multi-participante, deliberadamente sintético, que instancia el perfil sobre un escenario de predicción de riesgo y permite evaluar la adecuación representacional sin presentarlo como despliegue productivo ni validación externa por expertos.
\end{enumerate}

La evidencia reportada es, por tanto, evidencia ejecutable sobre el artefacto ontológico y sobre un grafo de demostración. No se interpreta como certificación legal, adopción institucional ni validación de una plataforma federada en producción. El resto del artículo revisa el trabajo relacionado, describe la metodología y los requisitos, presenta la ontología, reporta la evaluación, discute límites y rutas de evolución, y resume la disponibilidad del artefacto.

% =====================================================================
\section{Trabajo relacionado}\label{sec:related}

DAIMO se sitúa en la intersección de cuatro familias de trabajo que la literatura existente trata habitualmente por separado: documentación narrativa de modelos, ontologías del ciclo de vida de aprendizaje automático, vocabularios de gobernanza y estándares de espacios de datos. Revisamos cada familia por orden de proximidad al problema, identificando qué aporta al puente catálogo-dataspace, dónde lo deja abierto y cómo DAIMO se posiciona respecto a ella. Una quinta subsección revisa las ontologías SWJ recientes cuya metodología de validación hemos adoptado. La Tabla~\ref{tab:reuse-matrix}, inspirada en la matriz de reuso de RePlanIT~\cite{kurteva}, compara todas las familias contra los cinco atributos operacionales que un perfil para modelos de IA gobernados debe cubrir.

\subsection{Descripción de modelos de IA}

\paragraph{Documentación narrativa.}

Las \emph{model cards} de Mitchell et al.~\cite{mitchell19}, las \emph{factsheets} de Arnold et al.~\cite{arnold19} y las \emph{datasheets for datasets} de Gebru et al.~\cite{gebru21} elevaron los estándares de transparencia sobre artefactos de IA. Describen en lenguaje natural el uso previsto del modelo, sus limitaciones, el rendimiento esperado sobre benchmarks conocidos y las condiciones bajo las que ha sido entrenado y evaluado. Son valiosas para la comunicación con lectores humanos y para la rendición de cuentas cualitativa, pero insuficientes como sustrato para descubrimiento automatizado, filtrado por política o vinculación de evidencia en flujos máquina a máquina: un espacio de datos requiere primitivas que los agentes de software puedan consumir, no documentos que otros agentes deban interpretar.

\paragraph{Ontologías del ciclo de vida ML.}

Allí donde la documentación narrativa falla en el consumo máquina, las ontologías estructuradas del ciclo de vida ML ocupan su lugar. EXPO~\cite{expo} ofrece una ontología general de experimentos científicos que ha servido de base a trabajos posteriores. ML-Schema~\cite{mls} formaliza la semántica de experimentos de aprendizaje automático mediante un esquema reutilizable y ligero. MEX~\cite{mex} provee un vocabulario de intercambio para resultados de experimentos, y OntoDM~\cite{ontodm} junto con DMOP~\cite{dmop} articulan el dominio de la minería de datos y el \emph{semantic meta-mining}. Souza et al.~\cite{souza19} y Schmidt et al.~\cite{schmidt21} extienden la perspectiva hacia la procedencia y la FAIRness del ciclo de vida ML. Estas ontologías proporcionan un fundamento riguroso pero permanecen centradas en el experimento científico y en el proceso de entrenamiento; ninguna aborda la dimensión contractual del intercambio entre participantes autónomos.

Un grupo reciente de trabajos se desplaza desde el experimento hacia la taxonomía y gestión de activos de IA. La \emph{Artificial Intelligence Ontology} (AIO) sistematiza conceptos técnicos y éticos de IA en ramas como sesgo, capas, tareas de aprendizaje automático, funciones matemáticas, modelos, redes, preprocesamiento y estrategias de entrenamiento~\cite{aio}. La \emph{Intelligence Task Ontology} (ITO) organiza tareas, benchmarks y métricas de rendimiento para analizar el progreso de capacidades de IA a gran escala~\cite{ito}. En la misma línea, Novacek et al.\ proponen una gestión ontológica de modelos y datasets orientada a relacionar modelos con los datos usados para su entrenamiento y validación~\cite{novacek24}. Estos enfoques refuerzan la necesidad de describir la estructura técnica del modelo, sus datos y sus métricas; DAIMO aborda una pregunta distinta y complementaria: cómo convertir esas descripciones en un activo ofertable, autorizable, invocable, auditable y comparable dentro de un espacio de datos federado.

El trabajo más próximo al dominio de DAIMO, y el más reciente, es MLDCAT-AP~3.0.0~\cite{mldcatap}, publicado por SEMIC en 2025. MLDCAT-AP extiende DCAT para describir modelos de IA como subclases de \code{dcat:Dataset}, con metadatos ricos sobre política, benchmark, evaluación, ejecución, riesgo e infraestructura computacional. DAIMO lo reutiliza directamente: la clase \code{it6:MachineLearningModel} es la clase del modelo en DAIMO, y las clases \code{it6:Task}, \code{it6:Run}, \code{it6:Flow}, \code{it6:Evaluation} e \code{it6:ComputerInfrastructure} capturan respectivamente la tarea objetivo, la ejecución concreta, el \emph{pipeline} asociado, la medición resultante y el entorno de cómputo. MLDCAT-AP es fuerte en el lado del modelo pero deliberadamente agnóstico respecto al espacio de datos en el que el modelo circula: no modela el acto de ofrecer un modelo a través de fronteras organizativas, ni las nociones de despliegue gobernado, autorización de ejecución anclada a un acuerdo o procedencia agregada entre contextos administrativos distintos. MLDCAT-AP cubre el modelo; DAIMO cubre el puente entre el modelo y el dataspace.

La Tabla~\ref{tab:mldcatap-delta} explicita la relación clase a clase entre DAIMO y MLDCAT-AP. El patrón es intencional: DAIMO adopta las clases de MLDCAT-AP para el plano del activo de IA, añade restricciones de perfil cuando el intercambio gobernado exige completitud operativa, e introduce extensiones sólo para las relaciones puente que MLDCAT-AP no cubre. La versión evaluada del perfil se mantiene sobre MLDCAT-AP~3.0.0 porque los alineamientos, ejemplos y pruebas SHACL fueron construidos contra esa especificación. MLDCAT-AP~3.1.0, publicada como \emph{Candidate Recommendation} en 2026~\cite{mldcatap31}, queda como riesgo de evolución que deberá contrastarse en futuras versiones del alineamiento.

\begin{table*}[t]
\small\sf\centering
\caption{Relación entre DAIMO y MLDCAT-AP~3.0.0, clase por clase. \emph{Adoptada} indica que el constructo permanece en la capa MLDCAT-AP; \emph{perfilada}, que DAIMO añade expectativas de conformidad; \emph{extensión puente}, que DAIMO introduce un constructo para la relación catálogo-dataspace ausente en MLDCAT-AP.\label{tab:mldcatap-delta}}
\begin{tabular}{@{}p{4.4cm}p{2.5cm}p{7.1cm}@{}}
\toprule
\textbf{Clase MLDCAT-AP} & \textbf{Tratamiento DAIMO} & \textbf{Notas} \\
\midrule
\code{it6:MachineLearning}\newline\code{Model} & Adoptada + perfilada & Descripción del activo de IA; DAIMO exige política, título e identificador para intercambio gobernado. \\
\code{it6:Task} & Adoptada & Tarea objetivo del modelo. \\
\code{it6:Run} & Adoptada + perfilada & Registro de ejecución; DAIMO exige \emph{flow}, algoritmo, agente, marca temporal y enlace de autorización. \\
\code{it6:Flow} & Adoptada & Pipeline asociado a una ejecución. \\
\code{it6:Evaluation} & Adoptada & Resultado de medición; puede formar parte de la procedencia cruzada. \\
\code{it6:Evaluation}\newline\code{Measure} & Adoptada & Tipo y valor de métrica. \\
\code{it6:Benchmark} & Adoptada & Referencia de benchmark para condiciones de comparabilidad. \\
\code{it6:Computer}\newline\code{Infrastructure} & Adoptada & Entorno de cómputo donde se ejecuta un despliegue. \\
\code{it6:Hardware}, \code{it6:Library} & Adoptadas & Componentes del entorno de cómputo. \\
\code{it6:HarmRisk} & Adoptada & Anotación de riesgo sobre el modelo. \\
\code{it6:Modality} & Adoptada & Modalidad de entrada. \\
\code{it6:trainedOn}, \code{it6:testedOn} & Adoptadas & Enlaces a datasets de entrenamiento y prueba. \\
\midrule
\multicolumn{2}{@{}l}{\textbf{Constructos puente ausentes en MLDCAT-AP y añadidos por DAIMO}} & \\
\midrule
(sin análogo MLDCAT-AP) & Extensión puente & Oferta de catálogo gobernada con política ODRL asociada. \\
(sin análogo MLDCAT-AP) & Extensión puente & Autorización de ejecución basada en un acuerdo. \\
(sin análogo MLDCAT-AP) & Extensión puente & Relación de despliegue entre modelo, infraestructura y servicio invocable. \\
(sin análogo MLDCAT-AP) & Extensión puente & Formatos de entrada/salida y condiciones de autenticación. \\
(sin análogo MLDCAT-AP) & Extensión puente & Salida gobernable de una ejecución autorizada. \\
(sin análogo MLDCAT-AP) & Extensión puente & Evidencia de integridad y auditoría orientada a cumplimiento. \\
(sin análogo MLDCAT-AP) & Extensión puente & Agregación de procedencia entre contextos de participante. \\
(sin análogo MLDCAT-AP) & Extensión puente & Condiciones reificadas de comparabilidad de evaluaciones. \\
(sin análogo MLDCAT-AP) & Extensión puente & Rol de participante en el intercambio dentro del dataspace. \\
\bottomrule
\end{tabular}
\end{table*}

\subsection{Gobernanza de recursos y espacios de datos}

\paragraph{Vocabularios W3C de gobernanza.}

Las ontologías del ciclo de vida describen qué es un modelo y cómo se produjo, pero no gobiernan cómo se publica, se licencia o se atribuye. Tres vocabularios W3C proporcionan ese sustrato de gobernanza. DCAT~2~\cite{dcat2} proporciona las primitivas de catálogo sobre las que se apoyan tanto MLDCAT-AP como DAIMO. ODRL~2.2~\cite{odrl22} ofrece un modelo consolidado para expresar permisos, prohibiciones, obligaciones, ofertas (\code{odrl:Offer}) y acuerdos (\code{odrl:Agreement}) vinculantes; DAIMO lo usa intensivamente tanto para la política adherida a las ofertas como para los acuerdos que autorizan ejecuciones concretas. PROV-O~\cite{provo} aporta la ontología de procedencia con actividades, entidades, agentes, bundles y roles, y es la capa natural sobre la que codificar la trazabilidad de ejecuciones individuales y su agregación entre participantes. Pandit y Esteves~\cite{pandit-esteves} muestran cómo ODRL se combina con DPV para escenarios de salud, un patrón que DAIMO transpone al dominio IA. Estos vocabularios son maduros y ampliamente adoptados, pero por construcción son agnósticos al dominio: ninguno es específico de inteligencia artificial ni de espacios de datos.

\paragraph{Estándares de espacios de datos.}

Mientras que los vocabularios anteriores describen artefactos, los estándares de espacios de datos describen el intercambio en sí. La perspectiva desplaza el foco desde los artefactos aislados hacia los intercambios gobernados entre actores autónomos. Franklin et al.~\cite{franklin05} formularon las \emph{dataspaces} como respuesta a la heterogeneidad persistente de fuentes. Otto et al.~\cite{otto19} y Cappiello et al.~\cite{cappiello22} convierten la soberanía del dato, los contratos y la interoperabilidad controlada en supuestos explícitos del diseño. El \emph{Dataspace Protocol}~\cite{dsp} es la especificación W3C-CG que define los mensajes de catálogo, la negociación de contratos y los procesos de transferencia; Eclipse Dataspace Components~\cite{edc} es su implementación de referencia.

Conviene distinguir entre \emph{framework} y \emph{vocabulario}: EDC es un framework de ejecución construido en Java, no una ontología. La semántica que EDC implementa proviene de DSP (protocolo), ODRL (políticas) y DCAT (catálogo); sólo un puñado de conceptos son verdaderas extensiones específicas de EDC, y de ellos DAIMO referencia explícitamente \code{edc:ParticipantContext} por su utilidad para anclar la procedencia agregada a un contexto administrativo. Esta aclaración es importante porque en la literatura de espacios de datos las dos capas se confunden con frecuencia. La versión evaluada de DAIMO conserva los términos DSP utilizados durante el desarrollo; la evolución de DSP~2025-1~\cite{dsp2025} constituye un punto de seguimiento para futuras versiones del alineamiento informativo. Trabajos sobre ecosistemas federados como ConSolid~\cite{werbrouck24} muestran además que la Web Semántica evoluciona hacia escenarios multi-actor donde gobernanza e interoperabilidad deben diseñarse conjuntamente.

\subsection{Ingeniería ontológica ejecutable en SWJ}

Más allá del sustrato de dominio que las cuatro familias anteriores proporcionan, un quinto hilo de trabajo reciente informa cómo debe evaluarse DAIMO. El \emph{Semantic Web Journal} ha publicado recientemente un conjunto de ontologías que comparten un compromiso metodológico explícito con la validación ejecutable y la reproducibilidad. Kurteva et al.\ presentan RePlanIT~\cite{kurteva}, una ontología para pasaportes digitales de producto de equipos ICT, con SHACL y evaluación de usabilidad con expertos de industria. Palma et al.\ operacionalizan información de suelos a escala global en GloSIS~\cite{palma} mediante un perfil modular alineado con SOSA/SSN e ISO 28258. Woods et al.\ conectan ingeniería ontológica con procesamiento del lenguaje natural para actividades de mantenimiento industrial~\cite{woods}. Bareedu et al.\ derivan reglas SHACL desde estándares industriales OPC~UA~\cite{bareedu24}. Ferranti et al.\ formalizan restricciones de propiedad de Wikidata mediante SHACL y SPARQL~\cite{ferranti}. Robaldo y Batsakis~\cite{robaldo} analizan la interacción entre validación SHACL e inferencia sobre la Ontología del Tiempo. Penteado et al.~\cite{penteado23} estudian metodologías de publicación de linked open government data. Estos trabajos establecen colectivamente la expectativa ---que DAIMO adopta como vinculante--- de que una ontología con pretensiones operacionales presente evidencia reproducible de su corrección, no solamente una descripción axiomática.

\begin{table*}[t]
\small\sf\centering
\caption{Matriz de reuso: familias de trabajo relacionado contra los cinco atributos operacionales que un perfil para modelos de IA gobernados debe cubrir. $\bullet$ = cubierto nativamente; $\circ$ = cubierto parcialmente mediante extensiones; --- = no cubierto.\label{tab:reuse-matrix}}
\begin{tabular}{@{}lccccc@{}}
\toprule
\textbf{Familia / artefacto} & \textbf{Publicación} & \textbf{Política} & \textbf{Invocación} & \textbf{Trazabilidad} & \textbf{Evaluación} \\
\midrule
Model cards, factsheets, datasheets~\cite{mitchell19,arnold19,gebru21} & $\circ$ & --- & --- & --- & $\circ$ \\
ML-Schema, MEX, OntoDM, DMOP~\cite{mls,mex,ontodm,dmop} & --- & --- & --- & $\circ$ & $\bullet$ \\
AIO, ITO, gestión modelo-dataset~\cite{aio,ito,novacek24} & --- & --- & --- & $\circ$ & $\bullet$ \\
MLDCAT-AP~3.0.0~\cite{mldcatap} & $\bullet$ & $\circ$ & --- & $\bullet$ & $\bullet$ \\
DCAT~2~\cite{dcat2} & $\bullet$ & --- & $\circ$ & --- & --- \\
ODRL~2.2~\cite{odrl22} & --- & $\bullet$ & --- & --- & --- \\
PROV-O~\cite{provo} & --- & --- & --- & $\bullet$ & $\circ$ \\
DSP / EDC~\cite{dsp,edc} & $\circ$ & $\circ$ & $\bullet$ & --- & --- \\
\midrule
\textbf{DAIMO} (este trabajo) & $\bullet$ & $\bullet$ & $\bullet$ & $\bullet$ & $\bullet$ \\
\bottomrule
\end{tabular}
\end{table*}

La lectura horizontal de la Tabla~\ref{tab:reuse-matrix} indica que ningún artefacto individual cubre simultáneamente las cinco dimensiones; la lectura vertical confirma que cada columna está cubierta nativamente por al menos un vocabulario reutilizable. DAIMO se diseña, por tanto, como un perfil de integración que ata esas cinco capacidades a través de alineamientos axiomáticos verificados, sin duplicar lo que los vocabularios maduros ya resuelven.

% =====================================================================
\section{Metodología y requisitos}\label{sec:method}

DAIMO adopta la metodología \emph{Linked Open Terms} (LOT)~\cite{lot22} por dos razones. Primera, su orientación industrial: LOT prescribe un flujo ligero con entregables en cada una de sus cuatro fases (especificación de requisitos, implementación, publicación y mantenimiento). Segunda, LOT trata la reutilización como fase de primera clase, no como subproducto, lo que encaja con la regla reuse-first que el problema impone. Las subsecciones siguientes describen el escenario ejecutivo que guió la elicitación de requisitos (§\ref{sec:scenario}), las preguntas de competencia y los requisitos no funcionales resultantes (§\ref{sec:method-req}), y las tres decisiones centrales de diseño que dieron forma a la ontología (§\ref{sec:method-design}).

\subsection{Escenario ejecutivo}\label{sec:scenario}

La especificación de requisitos parte de un escenario ejecutivo anonimizado que posteriormente actúa como caso de estudio (§\ref{sec:eval-usecase}). Un Proveedor~A publica un modelo de predicción geoespacial de riesgo en un catálogo federado gestionado por una Plataforma~C. Un Consumidor municipal~B necesita invocar el modelo para apoyar decisiones operativas tras un evento extremo, pero no puede hacerlo sin antes descubrirlo bajo una política compatible con su uso público, verificar que satisface un umbral mínimo de calidad sobre un contexto de evaluación común, negociar un acuerdo que autorice sus ejecuciones concretas e invocarlo a través de un endpoint autenticado. Un Evaluador científico~D compara el modelo con una línea base sobre el mismo contexto de evaluación. Un Actor de cumplimiento~E audita las ejecuciones y archiva la evidencia correspondiente. Los nombres, modelos, endpoints, métricas y fechas del escenario son sintéticos por diseño: el caso no se presenta como piloto desplegado, sino como demostrador reproducible de adecuación representacional.

Este escenario sintetiza las cinco tareas operativas que una ontología del puente catálogo-dataspace debe soportar simultáneamente (publicación, descubrimiento, invocación, trazabilidad y evaluación) e involucra a cinco tipos de actor claramente diferenciados cuyas necesidades son complementarias pero no coinciden. Es verosímil: no depende de tecnología hipotética y se puede instanciar con los vocabularios reutilizados más las extensiones que DAIMO introduce.

\subsection{Especificación de requisitos}\label{sec:method-req}

Los requisitos se elicitaron desde tres fuentes complementarias: el escenario anterior, el dosier interno de MLDCAT-AP~3.0.0 y el análisis de las extensiones Eclipse EDC empleadas por INESData. A partir de ellas se formularon veintitrés preguntas de competencia, organizadas en cinco categorías por actor y etapa del ciclo de vida. La categoría~R (registro y publicación, cinco preguntas) cubre las necesidades del proveedor al publicar el modelo como activo gobernado con identidad, licencia, política y contrato de invocación. La categoría~D (descubrimiento y selección, cuatro preguntas) recoge las preguntas que un consumidor plantea para localizar un modelo adecuado por tarea, dominio, política y umbral de calidad. La categoría~E (ejecución y auditabilidad, cinco preguntas) agrupa las preguntas que un operador y un actor de gobernanza se plantean para invocar servicios, trazar ejecuciones y reconstruir evidencia. La categoría~V (evaluación y reproducibilidad, cinco preguntas) contiene las preguntas del evaluador para comparar resultados bajo condiciones homogéneas. Finalmente, la categoría~G (puente de gobernanza, cuatro preguntas) cubre las preguntas multi-actor cuya respuesta exige exclusivamente semántica nativa de DAIMO y que, como se argumenta en §\ref{sec:discussion}, no pueden responderse con los vocabularios reutilizados aisladamente. La Tabla~\ref{tab:actors} resume la correspondencia actor-categoría; la Tabla~\ref{tab:req-trace} traza cada familia de requisitos hacia los compromisos de modelado que la satisfacen; la Tabla~\ref{tab:domain-terms} normaliza los términos nucleares del dominio antes de introducir clases o propiedades; el conjunto completo de las veintitrés preguntas en lenguaje natural aparece en el Apéndice~\ref{app:cqs}.

\begin{table*}[t]
\small\sf\centering
\caption{Actores involucrados en el escenario ejecutivo y categorías de preguntas de competencia derivadas.\label{tab:actors}}
\begin{tabular}{@{}p{3.5cm}p{2cm}p{7cm}p{2.5cm}@{}}
\toprule
\textbf{Actor} & \textbf{Categoría} & \textbf{Necesidad} & \textbf{CQs} \\
\midrule
Proveedor de modelos & R & Publicar un modelo como activo gobernado con identidad, licencia, política y contrato de invocación & \cq{CQ-R1}--\cq{CQ-R5} \\
Consumidor de modelos & D & Descubrir modelos por tarea, dominio, política y umbral de calidad & \cq{CQ-D1}--\cq{CQ-D4} \\
Operador de plataforma & E & Invocar servicios, trazar ejecuciones y reconstruir evidencia operativa & \cq{CQ-E1}--\cq{CQ-E5} \\
Evaluador & V & Comparar resultados en un contexto compartido e inspeccionar reproducibilidad & \cq{CQ-V1}--\cq{CQ-V5} \\
Gobernanza (multi-actor) & G & Puentear oferta, despliegue, autorización y procedencia entre participantes & \cq{CQ-G1}--\cq{CQ-G4} \\
\bottomrule
\end{tabular}
\end{table*}

\begin{table*}[t]
\small\sf\centering
\caption{Trazabilidad desde requisito operacional hacia compromiso de modelado. La tabla explicita por qué existe cada grupo de modelado y qué artefacto de evaluación lo ejercita.\label{tab:req-trace}}
\begin{tabular}{@{}p{3.2cm}p{4.3cm}p{4.2cm}p{2.8cm}@{}}
\toprule
\textbf{Familia de requisito} & \textbf{Pregunta que el perfil debe responder} & \textbf{Compromiso DAIMO principal} & \textbf{Evidencia} \\
\midrule
Publicación & ¿Este modelo es un activo gobernado en un contexto de catálogo concreto? & Registro de oferta que vincula modelo, proveedor, contexto de catálogo y política aplicable. & CQs R1--R5, G1 \\
Descubrimiento & ¿Puede un consumidor filtrar modelos por tarea, dominio, política y umbral de evaluación? & Semántica de modelo, tarea y evaluación de MLDCAT-AP más enlaces a política y contexto de evaluación. & CQs D1--D4 \\
Invocación & ¿Qué servicio, método de autenticación y contrato E/S aplican a un modelo seleccionado? & Despliegue de modelo conectado a servicios DCAT y contratos E/S, permitiendo varios endpoints por modelo alojado. & CQs E1, G2 \\
Autorización & ¿Qué acuerdo autorizó una ejecución concreta por un participante concreto? & Acuerdo negociado enlazado al agente beneficiario y a la ejecución que cubre. & CQs E2, G3; INV-2, INV-4 \\
Auditabilidad & ¿Puede un actor de gobernanza reconstruir artefactos, evidencias y procedencia entre contextos? & Artefacto derivado, evidencia de auditoría y bundle de procedencia entre participantes. & CQs E4, E5, G4; INV-1 \\
Evaluación & ¿Son comparables las mediciones bajo la misma tarea, versión de dataset, protocolo y semilla? & Contexto compartido de evaluación que fija tarea, dataset, versión, protocolo y semilla. & CQs V1--V5 \\
\bottomrule
\end{tabular}
\end{table*}

El flujo operacional de ejecución federada se interpreta en DAIMO como una cadena semántica auditable, no como una arquitectura de runtime. El descubrimiento se representa mediante ofertas gobernadas y CQs D1--D4; la negociación, mediante \code{odrl:Agreement} y \code{ExecutionAuthorization}; el despliegue, mediante \code{ModelDeployment}, \code{dcat:DataService} e \code{IOContract}; la inferencia, mediante \code{it6:Run} y \code{DerivedArtifact}; y el cierre operativo o \emph{clearing}, mediante \code{AuditEvidence} y \code{CrossParticipantProvenanceRecord}. Esta lectura preserva la premisa de soberanía: DAIMO puede describir ejecuciones en el nodo del proveedor de datos, en un nodo neutral o en una infraestructura acordada, pero no prescribe si el modelo viaja al dato, el dato al modelo, o si la ejecución ocurre dentro de un mecanismo particular de aislamiento.

La enumeración de términos se trató como un artefacto intermedio entre las preguntas de competencia y la TBox. La Tabla~\ref{tab:domain-terms} no añade semántica nueva: normaliza el vocabulario operativo que aparece en el escenario, decide si cada término se reutiliza o se especializa, y evita convertir temas generales en clases nativas sin necesidad. Las definiciones formales completas se mantienen en la ontología mediante \code{skos:definition}, \code{skos:example}, dominios, rangos y shapes.

\begin{table*}[t]
\scriptsize\sf\centering
\caption{Glosario inicial de términos nucleares del dominio y su realización ontológica.\label{tab:domain-terms}}
\begin{tabular}{@{}p{2.8cm}p{4.9cm}p{6.8cm}@{}}
\toprule
\textbf{Término del dominio} & \textbf{Sentido operacional} & \textbf{Realización ontológica} \\
\midrule
Modelo de IA & Artefacto entrenado o empaquetado que puede publicarse, evaluarse e invocarse. & Reutilizado mediante la clase de modelo de MLDCAT-AP; DAIMO no duplica la clase del modelo. \\
Oferta gobernada & Acto de registrar un modelo como activo disponible bajo política en un catálogo federado. & \code{AIAssetOffering} enlaza modelo, agente y \code{odrl:Offer}; se alinea con \code{dcat:CatalogRecord}. \\
Participante & Agente jurídico u organizativo que actúa dentro del espacio de datos. & \code{foaf:Agent} para el agente; \code{ParticipantRole} para el rol contextual. \\
Rol de participante & Papel asumido respecto al activo: proveedor, consumidor, operador, evaluador o gobernanza. & Cinco subclases no disjuntas de \code{ParticipantRole}, porque un agente puede asumir varios roles. \\
Contexto de participante & Ámbito administrativo donde se publica, negocia, opera o audita el intercambio. & Reutilizado como \code{edc:ParticipantContext}; DAIMO lo usa para acotar roles y procedencia. \\
Política y acuerdo & Condiciones de uso de la oferta y resultado de una negociación de acceso. & Reutilización de \code{odrl:Offer} y \code{odrl:Agreement}; DAIMO especializa sólo la autorización de ejecución. \\
Despliegue de modelo & Instancia alojada que sirve un modelo mediante uno o más servicios. & \code{ModelDeployment} conecta modelo, infraestructura, servicio DCAT y contrato E/S. \\
Servicio o endpoint & Interfaz accesible para invocar un despliegue. & Reutilizado como \code{dcat:DataService}; DAIMO añade el enlace desde el despliegue. \\
Contrato E/S & Formato de entrada, formato de salida, autenticación y esquemas opcionales. & \code{IOContract} con \code{inputFormat}, \code{outputFormat}, \code{authMethod} y esquemas opcionales. \\
Ejecución autorizada & Invocación concreta cubierta por un acuerdo válido. & \code{it6:Run} para la ejecución y \code{ExecutionAuthorization} para el acuerdo que la autoriza. \\
Artefacto derivado & Resultado gobernable producido por una ejecución. & \code{DerivedArtifact}, enlazado a la ejecución y a la autorización bajo la que se produjo. \\
Evidencia de auditoría & Registro verificable que respalda una actividad ejecutada. & \code{AuditEvidence} con actividad, firmante, marca temporal y checksum SPDX. \\
Procedencia cruzada & Narrativa de procedencia que cruza al menos dos contextos de participante. & Clase DAIMO de procedencia cruzada como bundle PROV con contextos explícitos. \\
Contexto de evaluación & Condiciones que hacen comparables las métricas de dos modelos. & \code{SharedEvaluationContext} fija tarea, dataset, versión, protocolo, semilla y flow. \\
Evaluación y métrica & Medición de desempeño usada para comparar modelos. & Reutilización de \code{it6:Evaluation}; DAIMO enlaza la evaluación al contexto compartido. \\
\bottomrule
\end{tabular}
\end{table*}

Once de las veintitrés preguntas exigen razonamiento explícito (rutas de subclase, cadenas de subpropiedad o agregación SPARQL), mientras que las doce restantes son consultas directas sobre datos DAIMO-nativos. Esta proporción es deliberada: un perfil que se apoyara exclusivamente en consultas directas no ejercitaría la parte del artefacto que hace valiosa la semántica OWL, y un perfil que exigiera razonamiento para toda pregunta sería impracticable en operación sobre endpoints reales. Los requisitos no funcionales complementan los funcionales y fijan: perfil OWL~2~DL (para mantener decidibilidad y compatibilidad con HermiT), licencia CC-BY~4.0 para la ontología y Apache~2.0 para el código de validación, identificadores persistentes bajo \url{https://w3id.org/pionera/daimo}, serializaciones RDF estándar y validación SHACL accesible.

\subsection{Decisiones centrales de diseño}\label{sec:method-design}

Tres decisiones estructurales se tomaron en la fase de conceptualización y condicionan el resto del diseño.

\paragraph{Decisión~1: reutilización preferente.} DAIMO reutiliza directamente las clases de MLDCAT-AP para todo aquello que MLDCAT-AP ya modela con precisión. \code{it6:MachineLearningModel} es la clase del modelo; \code{it6:ComputerInfrastructure} describe el entorno de ejecución; \code{it6:Run}, \code{it6:Flow} e \code{it6:Evaluation} capturan respectivamente la ejecución concreta, el \emph{pipeline} asociado y la medición resultante; \code{dcat:Dataset}, \code{dcat:Distribution} y \code{dcat:DataService} publican el activo; \code{odrl:Offer} y \code{odrl:Agreement} expresan la política y el acuerdo. Introducir equivalentes DAIMO-nativos de cualquiera de estas clases violaría el principio de compromiso ontológico mínimo de Gruber y fracturaría el ecosistema DCAT-AP sin ganancia; una clase DAIMO-nativa sólo se justifica cuando la unión de los vocabularios reutilizados deja un hueco concreto que impide responder a al menos una pregunta de competencia.

Esta decisión ordena el modelo conceptual sin convertir cada tema en una clase nueva. El modelo de IA se mantiene como \code{it6:MachineLearningModel}; los recursos de datos como \code{dcat:Dataset} y \code{dcat:Distribution}; la infraestructura como \code{it6:ComputerInfrastructure}; las métricas como \code{it6:Evaluation} y \code{it6:EvaluationMeasure}; y la gobernanza como combinación de políticas ODRL, autorizaciones DAIMO, roles de participante y evidencia auditable. El linaje no se reifica como una clase monolítica, sino como una cadena consultable que conecta datos, ejecución, despliegue, artefactos derivados y bundles PROV. Por la misma razón, candidatos habituales como \code{AIModel}, \code{DataSet}, \code{Algorithm}, \code{ModelVersion}, \code{hasOwner}, \code{hasLicense}, frameworks o hiperparámetros se reutilizan desde MLDCAT-AP, DCAT, ODRL, PROV-O o ML-Schema, o quedan como metadatos de perfiles externos si ninguna pregunta de competencia del núcleo los exige. Así, DAIMO permanece desacoplada de frameworks concretos como PyTorch o TensorFlow y evita importar modelos IDS o Gaia-X como dependencias del núcleo; esos ecosistemas quedan como contextos de adopción interoperables mediante DCAT, ODRL, PROV-O, DSP y los alineamientos del perfil.

\paragraph{Decisión~2: separación en tres capas complementarias.} La arquitectura distingue tres capas que DAIMO conecta pero no sustituye. La capa~A es el plano del espacio de datos: DSP con las extensiones EDC puntuales que DAIMO referencia. La capa~B es el plano de gobernanza de plataforma: INESData (o una plataforma análoga) con sus mecanismos de registro federado y políticas operativas. La capa~C es el plano del activo de IA/ML: MLDCAT-AP y los vocabularios que lo soportan. DAIMO es el perfil de integración que conecta A, B y C; no es una cuarta capa paralela que replique funcionalidad, sino un puente que hace consultable la articulación entre las tres.

\paragraph{Decisión~3: una clase nativa por hueco, no por tema.} Dada la reutilización preferente, cada clase DAIMO-nativa debe justificar su existencia por un hueco concreto en la unión de los vocabularios reutilizados y por al menos una pregunta de competencia que no puede responderse sin ella. El criterio se aplica con rigor y produce nueve clases de primer nivel (\code{AIAssetOffering}, \code{ParticipantRole}, \code{ModelDeployment}, \code{IOContract}, \code{ExecutionAuthorization}, \code{DerivedArtifact}, \code{CrossParticipantProvenanceRecord}, \code{AuditEvidence} y \code{SharedEvaluationContext}) más cinco subclases de rol. Cualquier clase que no cumpla simultáneamente ambas condiciones queda descartada.

% =====================================================================
\section{La ontología DAIMO}\label{sec:ontology}

\subsection{Arquitectura modular y métricas}\label{sec:scope}

DAIMO modela el ciclo de vida operacional de un activo de modelo de IA dentro de un espacio de datos gobernado: publicación como registro de catálogo, descubrimiento sensible a política y a calidad, invocación controlada mediante servicios tipados, trazabilidad de ejecuciones individuales y evaluación contextualizada y reproducible. El alcance se restringe deliberadamente a esta cadena operacional. Las responsabilidades ortogonales, como la modelización detallada del \emph{pipeline} de entrenamiento, la semántica fina de políticas de privacidad o la teoría formal de la derivación PROV, ya están cubiertas por vocabularios consolidados cuya duplicación fragmentaría el ecosistema. DAIMO aporta, en cambio, los constructos necesarios para que esos vocabularios se articulen coherentemente en el caso concreto de un modelo de IA que circula entre participantes autónomos.

La ontología se estructura en tres módulos complementarios, distribuidos en archivos Turtle separados para que un consumidor pueda cargar únicamente los que necesite. El \textbf{módulo núcleo} (\code{daimo-core.ttl}) contiene las catorce clases y las veintinueve propiedades de objeto más ocho propiedades de datatype DAIMO-nativas, los axiomas globales de disyunción y las caracterizaciones sobre propiedades (funcionales, asimétricas, inversas) que capturan las invariantes de tipo a nivel TBox. El \textbf{módulo de alineamientos} (\code{alignment.ttl}) encapsula ocho declaraciones \code{rdfs:subClassOf} externas y cinco declaraciones \code{rdfs:subPropertyOf} externas que conectan DAIMO con los vocabularios reutilizados, más las declaraciones mínimas de los términos externos que DAIMO referencia, cada una con un \code{rdfs:seeAlso} a la sección correspondiente de la especificación fuente; en conjunto, el paquete contiene seis subpropiedades DAIMO hacia vocabularios externos porque \code{hasOffering}~$\sqsubseteq$~\code{foaf:isPrimaryTopicOf} se declara en el núcleo como inversa de \code{offersModel}. El \textbf{módulo SHACL} (\code{daimo-shapes.ttl}) aporta las dieciocho \code{sh:NodeShape} del perfil: nueve estructurales, tres de conformidad sobre clases reutilizadas y seis invariantes entre clases mediante \code{sh:SPARQLConstraint}. El módulo se declara como \code{owl:Ontology} independiente con su propia \code{owl:versionIRI}, de modo que la capa de validación se libera y referencia como un artefacto con ciclo de vida propio.

La separación en tres módulos no es accidental. El núcleo puede ser consumido por una aplicación que sólo necesite las primitivas DAIMO; el módulo de alineamientos se carga cuando se quiere que un razonador infiera las relaciones entre DAIMO y los vocabularios reutilizados; el módulo SHACL se carga cuando se quiere validar un grafo concreto contra el perfil. La Figura~\ref{fig:architecture} ilustra la arquitectura modular y las relaciones con los vocabularios reutilizados.

\begin{table*}[t]
\scriptsize\sf\centering
\caption{Perfil técnico de implementación de DAIMO. La tabla traduce requisitos habituales de implementación a artefactos existentes del paquete, sin introducir módulos ni clases adicionales.\label{tab:implementation-profile}}
\begin{tabular}{@{}p{3.0cm}p{7.0cm}p{4.8cm}@{}}
\toprule
\textbf{Pilar} & \textbf{Implementación en DAIMO} & \textbf{Decisión de alcance} \\
\midrule
Lenguaje, serialización e IRIs & Núcleo OWL~2~DL en Turtle; serializaciones RDF/XML, JSON-LD y N-Triples preparadas para publicación; namespace canónico \code{https://w3id.org/pionera/daimo}. & La reproducibilidad local no depende de un endpoint persistente activo. \\
Vistas conceptuales de implementación & Metadatos de modelo: MLDCAT-AP; datos y catálogo: DCAT; ejecución e infraestructura: MLDCAT-AP más \code{ModelDeployment}/\code{IOContract}; gobernanza y confianza: ODRL, autorizaciones, evaluación contextualizada y evidencia. & Estas son vistas funcionales, no módulos ontológicos nuevos como \code{ModelMetadata} o \code{GovernanceTrust}. \\
Relaciones críticas & Se reutiliza \code{it6:trainedOn}; los formatos y esquemas se expresan en \code{IOContract}; el despliegue se conecta al endpoint mediante \code{exposedAs} y \code{dcat:DataService}; la conformidad normativa puede expresarse con \code{dct:conformsTo} o perfiles ODRL/SHACL específicos. & No se añaden propiedades duplicadas como \code{trainedOn}, \code{deployedIn} o \code{compliesWith} si ya existe una ruta estándar o perfilable. \\
Validación y razonamiento & SHACL estructural, SHACL sobre clases reutilizadas, seis invariantes SHACL-SPARQL, HermiT, OWL-RL, OOPS!, pruebas negativas y 23 consultas SPARQL. & No se exige una métrica universal como \emph{accuracy}; las métricas obligatorias dependen del contexto de evaluación y de la CQ ejercitada. \\
Almacenamiento y acceso & El paquete RDF puede cargarse en triplestores estándar y exponerse por SPARQL; el conjunto de consultas de competencia demuestra el patrón de acceso. & El artículo no afirma un endpoint productivo GraphDB/Fuseki; sólo compatibilidad técnica con ese tipo de despliegue. \\
\bottomrule
\end{tabular}
\end{table*}

\begin{figure}[t]
\centering
\begin{tikzpicture}[node distance=0.9cm and 0.7cm]
  \node[modulebox, fill=yellow!25] (core) {N\'ucleo DAIMO\\\scriptsize\normalfont\code{daimo-core.ttl}};
  \node[modulebox, fill=orange!20, below=of core] (align) {M\'odulo alineamientos\\\scriptsize\normalfont\code{alignment.ttl}};
  \node[modulebox, fill=red!15, below=of align] (shapes) {M\'odulo SHACL\\\scriptsize\normalfont\code{daimo-shapes.ttl}};

  \node[vocabbox, right=3.2cm of core, yshift=1.0cm] (mldcat) {MLDCAT-AP 3.0.0};
  \node[vocabbox, below=0.15cm of mldcat] (dcat) {DCAT 2};
  \node[vocabbox, below=0.15cm of dcat] (odrl) {ODRL 2.2};
  \node[vocabbox, below=0.15cm of odrl] (prov) {PROV-O};
  \node[vocabbox, below=0.15cm of prov] (dsp) {DSP / EDC};

  \draw[modulearrow] (align.east) -- ++(0.4,0) |- (mldcat.west);
  \draw[modulearrow] (align.east) -- ++(0.4,0) |- (dcat.west);
  \draw[modulearrow] (align.east) -- ++(0.4,0) |- (odrl.west);
  \draw[modulearrow] (align.east) -- ++(0.4,0) |- (prov.west);
  \draw[modulearrow] (align.east) -- ++(0.4,0) |- (dsp.west);

  \draw[objpropertyarrow] (core.south) -- node[right, font=\scriptsize\sffamily] {\code{owl:imports}} (align.north);
  \draw[objpropertyarrow] (shapes.north) -- node[right, font=\scriptsize\sffamily] {valida} (align.south);

  \begin{scope}[on background layer]
    \node[fit=(core)(shapes), draw=gray!50, dashed, inner sep=8pt, rounded corners] (daimobox) {};
    \node[above=0.0cm of daimobox, font=\sffamily\footnotesize\bfseries] {Perfil DAIMO};
  \end{scope}
  \begin{scope}[on background layer]
    \node[fit=(mldcat)(dsp), draw=gray!50, dashed, inner sep=8pt, rounded corners] (vocbox) {};
    \node[above=0.0cm of vocbox, font=\sffamily\footnotesize\bfseries] {Vocabularios reutilizados};
  \end{scope}
\end{tikzpicture}
\caption{Arquitectura modular de DAIMO y vocabularios reutilizados. Las cajas amarilla, naranja y roja representan los tres m\'odulos de DAIMO (n\'ucleo, alineamientos, shapes); las verdes, los vocabularios reutilizados. Las flechas discontinuas desde el m\'odulo de alineamientos portan las declaraciones \code{rdfs:subClassOf} y \code{rdfs:subPropertyOf}. El n\'ucleo importa alineamientos para inferencia; el m\'odulo SHACL valida datos contra el TBox combinado.}
\label{fig:architecture}
\end{figure}

\paragraph{Métricas del artefacto.}\label{sec:metrics}

Antes de recorrer las clases individuales que pueblan los tres módulos de §\ref{sec:scope}, reportamos el tamaño y la expresividad del artefacto globalmente. La Tabla~\ref{tab:metrics} extrae las métricas directamente de los archivos Turtle y de los reportes producidos por la suite de validación (§\ref{sec:eval}). Los números permiten a un lector interesado estimar el peso del artefacto antes de consumirlo y compararlo con ontologías publicadas previamente en la misma familia.

\begin{table}[t]
\small\sf\centering
\caption{Métricas del artefacto DAIMO verificadas contra los archivos del paquete.\label{tab:metrics}}
\begin{tabular}{@{}p{4.6cm}r@{}}
\toprule
\textbf{Dimensión} & \textbf{Valor} \\
\midrule
Perfil lógico declarado & OWL 2 DL \\
Clases nativas de primer nivel & 9 \\
Subclases nativas (\code{ParticipantRole}) & 5 \\
Total clases DAIMO-nativas & 14 \\
Propiedades de objeto (\code{owl:ObjectProperty}) & 29 \\
Propiedades de datatype (\code{owl:DatatypeProperty}) & 8 \\
Propiedades funcionales (\code{owl:FunctionalProperty}) & 27 \\
Propiedades asimétricas (\code{owl:AsymmetricProperty}) & 5 \\
Pares inversos explícitos (\code{owl:inverseOf}) & 5 \\
\code{rdfs:subClassOf} a externos & 8 \\
\code{rdfs:subPropertyOf} a externos & 6 \\
Axioma de disyunción nombrado & 1 \\
\code{sh:NodeShape} totales & 18 \\
\quad de completitud (una por clase DAIMO) & 9 \\
\quad de conformidad sobre reutilizadas & 3 \\
\quad invariantes \code{sh:SPARQLConstraint} & 6 \\
Preguntas de competencia & 23 \\
Pitfalls OOPS! (Crítico/Importante/Menor) & 0 / 0 / 2 \\
\bottomrule
\end{tabular}
\end{table}

\begin{table}[t]
\small\sf\centering
\caption{Desglose estilo OntoMetrics extraído de los módulos fuente DAIMO. Los conteos se restringen a sujetos del namespace DAIMO para evitar inflar la métrica con declaraciones externas usadas sólo para alineamiento.\label{tab:ontometrics}}
\begin{tabular}{@{}p{5.2cm}r@{}}
\toprule
\textbf{Tipo de axioma/declaración} & \textbf{Conteo} \\
\midrule
\code{owl:Class} & 14 \\
\code{owl:ObjectProperty} & 29 \\
\code{owl:DatatypeProperty} & 8 \\
\code{rdfs:subClassOf} & 13 \\
\code{rdfs:subPropertyOf} & 6 \\
\code{rdfs:domain} & 37 \\
\code{rdfs:range} & 37 \\
\code{owl:FunctionalProperty} & 27 \\
\code{owl:AsymmetricProperty} & 5 \\
\code{owl:inverseOf} & 5 \\
\code{owl:AllDisjointClasses} & 1 \\
\code{owl:equivalentClass} & 0 \\
\code{owl:disjointWith} explícitos & 0 \\
\bottomrule
\end{tabular}
\end{table}

\subsection{Clases nativas y alineamientos}\label{sec:classes}

Las nueve clases nativas de primer nivel, más las cinco subclases de \code{ParticipantRole}, constituyen el conjunto mínimo que las tres decisiones centrales de diseño producen. La Tabla~\ref{tab:classes} sintetiza cada una junto con su alineamiento axiomático a los vocabularios reutilizados; las subsecciones siguientes detallan las de mayor carga semántica. La Figura~\ref{fig:classes} muestra la vista general de clases y propiedades nativas.

\begin{table*}[t]
\small\sf\centering
\caption{Clases nativas de DAIMO, alineamientos axiomáticos y preguntas de competencia que responden.\label{tab:classes}}
\begin{tabular}{@{}p{4cm}p{3cm}p{6cm}p{2.2cm}@{}}
\toprule
\textbf{Clase} & \textbf{Alineamiento} & \textbf{Papel} & \textbf{CQs} \\
\midrule
\code{daimo:AIAsset}\newline\code{Offering} & \code{dcat:CatalogRecord} & Registro de catálogo que publica un modelo como oferta gobernada en un dataspace; porta política ODRL mediante \code{daimo:hasOfferPolicy}. & R1, R2, R5, G1 \\
\code{daimo:ParticipantRole} \newline (+5 subclases no disjuntas) & \code{prov:Role} & Rol asumido por un participante respecto a un activo de IA dentro de un contexto de dataspace. & R5 \\
\code{daimo:Model}\newline\code{Deployment} & \code{prov:Entity} & Instancia en ejecución o alojada del modelo, expuesta como \code{dcat:DataService} sobre una \code{it6:ComputerInfrastructure}. & E1, E3, G2 \\
\code{daimo:IOContract} & (sin alineamiento) & Contrato mínimo accionable: formatos de E/S, método de autenticación, esquemas opcionales. & R4, E1, G2 \\
\code{daimo:Execution}\newline\code{Authorization} & \code{odrl:Agreement} & Acuerdo ODRL derivado de una negociación DSP que autoriza ejecuciones concretas vía \code{daimo:authorizesRun}. & E2, E5, G3 \\
\code{daimo:DerivedArtifact} & \code{prov:Entity}, \code{dcat:Resource} & Salida gobernada y describible en catálogo producida por una ejecución bajo autorización. & E5, G4 \\
\code{daimo:CrossParticipant}\newline\code{ProvenanceRecord} & \code{prov:Bundle} & Agregación de actividades y entidades a través de contextos de participante para formar una narrativa de auditoría. & G4 \\
\code{daimo:AuditEvidence} & \code{prov:Entity} & Registro de evidencia de cumplimiento con \code{spdx:Checksum} estructurado, firmante y marca temporal. & E4 \\
\code{daimo:SharedEvaluation}\newline\code{Context} & (sin alineamiento) & Reificación de las condiciones de comparabilidad: tarea, dataset, versión, protocolo y semilla aleatoria. & V1, V2, V3 \\
\bottomrule
\end{tabular}
\end{table*}

\begin{figure*}[t]
\centering
\begin{tikzpicture}[
    x=0.76cm, y=0.95cm,
    daimoclass/.style={rectangle, draw=black!70, fill=yellow!25, thick,
      rounded corners=3pt, minimum width=2.3cm, minimum height=0.8cm,
      align=center, font=\ttfamily\scriptsize, inner sep=2pt},
    externclass/.style={rectangle, draw=black!55, fill=blue!12, thick,
      rounded corners=3pt, minimum width=2.1cm, minimum height=0.65cm,
      align=center, font=\ttfamily\scriptsize, inner sep=2pt},
    subarrow/.style={-{Triangle[open,length=3mm,width=2.8mm]},
      draw=blue!60!black, line width=0.9pt},
    proparrow/.style={-{Latex[length=2.2mm]}, draw=black!85, line width=0.6pt},
    edgelab/.style={fill=white, font=\scriptsize\sffamily, inner sep=1.5pt,
      text=black!85}
  ]

  \node[externclass] (catrec)    at ( 0,   5) {dcat:\\CatalogRecord};
  \node[externclass] (agreement) at ( 3.7, 5) {odrl:Agreement};
  \node[externclass] (role)      at ( 7,   5) {prov:Role};
  \node[externclass] (entity)    at (10.3, 5) {prov:Entity};
  \node[externclass] (resource)  at (13.4, 5) {dcat:Resource};
  \node[externclass] (bundle)    at (16.5, 5) {prov:Bundle};

  \node[daimoclass] (offering) at ( 0,   2.6) {daimo:\\AIAssetOffering};
  \node[daimoclass] (auth)     at ( 3.7, 2.6) {daimo:\\ExecutionAuthorization};
  \node[daimoclass] (role_d)   at ( 7,   2.6) {daimo:\\ParticipantRole};
  \node[daimoclass] (deploy)   at (10.3, 2.6) {daimo:\\ModelDeployment};
  \node[daimoclass] (derived)  at (13.4, 2.6) {daimo:\\DerivedArtifact};
  \node[daimoclass] (bundle_d) at (16.5, 2.6) {daimo:CrossPart.\\ProvenanceRecord};

  \node[daimoclass] (iocontract) at (10.3, 0.3) {daimo:IOContract};
  \node[daimoclass] (audit)      at (13.4, 0.3) {daimo:\\AuditEvidence};
  \node[daimoclass] (context)    at (16.5, 0.3) {daimo:Shared\\EvaluationContext};

  \draw[subarrow] (offering)     -- (catrec);
  \draw[subarrow] (auth)         -- (agreement);
  \draw[subarrow] (role_d)       -- (role);
  \draw[subarrow] (deploy)       -- (entity);
  \draw[subarrow] (derived.north) to[bend right=8] (entity.south);
  \draw[subarrow] (derived.north) to[bend left=8] (resource.south);
  \draw[subarrow] (bundle_d)     -- (bundle);
  \draw[subarrow] (audit.west)   to[bend left=35] (entity.south west);

  \draw[proparrow] (offering.south) to[bend right=32]
    node[edgelab, below, pos=0.5] {offeredBy}
    (role_d.south west);
  \draw[proparrow] (auth.east) --
    node[edgelab, midway] {grantedTo}
    (role_d.west);
  \draw[proparrow] (derived.north west) to[bend left=22]
    node[edgelab, above, pos=0.2, sloped] {underAuthorization}
    (auth.north east);
  \draw[proparrow] (deploy.south) --
    node[edgelab, right] {hasIOContract}
    (iocontract.north);
  \draw[proparrow] (derived.south) --
    node[edgelab, right] {hasAuditEvidence}
    (audit.north);
  \draw[proparrow] (bundle_d.south) to[bend left=15]
    node[edgelab, above, pos=0.4, sloped] {records}
    (audit.north east);

  \begin{scope}[shift={(0,-1.7)}]
    \node[daimoclass, minimum width=1.2cm, minimum height=0.45cm, inner sep=1pt]
      (lg1) at (0,0) {daimo:X};
    \node[anchor=west, font=\footnotesize\sffamily] at (1.0,0)
      {Clase DAIMO-nativa};
    \node[externclass, minimum width=1.2cm, minimum height=0.45cm, inner sep=1pt]
      at (5.2,0) {ext:Y};
    \node[anchor=west, font=\footnotesize\sffamily] at (6.1,0)
      {Clase reutilizada (externa)};
    \draw[subarrow] (10.3, 0) -- ++(0.8,0);
    \node[anchor=west, font=\footnotesize\sffamily] at (11.2,0)
      {\code{rdfs:subClassOf}};
    \draw[proparrow] (14.0, 0) -- ++(0.8,0);
    \node[anchor=west, font=\footnotesize\sffamily] at (14.9,0)
      {propiedad de objeto};
  \end{scope}

\end{tikzpicture}
\caption{Vista general de las clases nativas de DAIMO (amarillo) con sus alineamientos axiom\'aticos hacia los vocabularios reutilizados (azul) mediante \code{rdfs:subClassOf} (flechas triangulares huecas), y las principales propiedades nativas que articulan el perfil (flechas etiquetadas). \code{DerivedArtifact} tiene dos padres externos, \code{prov:Entity} y \code{dcat:Resource}; \code{AuditEvidence} aparece en la fila inferior pero es tambi\'en una \code{prov:Entity}. Las clases \code{IOContract} y \code{SharedEvaluationContext} carecen de alineamiento externo por dise\~no. Las propiedades cuyos destinos son clases reutilizadas (\code{offersModel} hacia \code{it6:MachineLearningModel}, \code{authorizesRun} hacia \code{it6:Run}) aparecen en la Tabla~\ref{tab:classes} y se omiten aqu\'i para no sobrecargar el diagrama.}
\label{fig:classes}
\end{figure*}

Las dos clases sin alineamiento externo, \code{IOContract} y \code{SharedEvaluationContext}, carecen de él deliberadamente. \code{IOContract} no es un \code{dcat:DataService} (un servicio es un recurso accesible, una URL, un endpoint, mientras que el contrato describe expectativas sintácticas y de autenticación sobre el intercambio) ni un \code{dct:Standard} (el contrato enumera obligaciones concretas, no refiere a una especificación externa). \code{SharedEvaluationContext} es una reificación metodológica, no una tarea: captura la combinación exacta bajo la que dos evaluaciones son comparables, no la tarea abstracta que dichas evaluaciones instancian. Alinearla a \code{owl:Thing} o \code{prov:Entity} añadiría ruido sin aportar semántica.

En el escenario ejecutivo de §\ref{sec:scenario}, el proveedor publica sus modelos mediante tres instancias de \code{AIAssetOffering} en el catálogo federado; cada oferta enlaza al modelo con \code{daimo:offersModel}, al proveedor con \code{daimo:offeredBy} y a la política ODRL con \code{daimo:hasOfferPolicy}. El despliegue del modelo seleccionado por el consumidor municipal se modela como una \code{ModelDeployment} que, mediante \code{daimo:exposedAs}, apunta a dos \code{dcat:DataService}s paralelos (REST y gRPC), cada uno con su \code{IOContract}. La autorización que el consumidor obtiene del proveedor tras la negociación DSP se reifica como \code{ExecutionAuthorization} y vincula, vía \code{daimo:authorizesRun}, las ejecuciones concretas que se permiten. La predicción producida por cada ejecución se reifica como \code{DerivedArtifact} con su \code{AuditEvidence} asociada, y el \code{CrossParticipantProvenanceRecord} agrega la cadena completa a través de los \code{edc:ParticipantContext}s del proveedor, el consumidor y la plataforma federada.

\medskip\noindent\emph{Registro y participación.}
\code{AIAssetOffering} y \code{ParticipantRole} resuelven conjuntamente el hueco del lado de publicación. MLDCAT-AP describe el modelo, y DCAT registra metadatos de catálogo, pero ninguno expresa el acto gobernado de poner un modelo a disposición de otro participante de un espacio de datos. La oferta se modela por tanto como el registro de catálogo de un modelo en un dataspace concreto, acotado por política, proveedor, modelo objetivo y contexto de participante. \code{ParticipantRole}, con cinco especializaciones no disjuntas, registra el rol que un agente jurídico asume respecto de un activo y un contexto de datos. El grupo responde a CQ-R1, CQ-R2, CQ-R5 y CQ-G1.

\medskip\noindent\emph{Plano de ejecución.}
\code{ModelDeployment} e \code{IOContract} cubren el hueco de invocación. DCAT proporciona servicios visibles desde catálogo, pero no especifica qué modelo aloja un servicio ni qué expectativas sintácticas y de autenticación debe satisfacer el consumidor para invocarlo. Un despliegue representa la instancia alojada o en ejecución de un modelo sobre infraestructura de cómputo. Puede exponerse mediante uno o varios servicios, y cada servicio puede declarar un contrato de invocación con formato de entrada/salida, método de autenticación y referencias opcionales a esquemas. La invocación no se modela como una propiedad directa modelo-consumidor: se reconstruye recorriendo oferta, despliegue, servicio, contrato E/S, autorización y \code{it6:Run}, lo que evita introducir relaciones redundantes como \code{isInvokedBy}. El perfil registra dónde y bajo qué contrato se puede invocar el modelo; no modela la orquestación de contenedores, el motor de enforcement ni el mecanismo TEE, que deben declararse como metadatos de infraestructura o perfiles específicos cuando una política los exige. El grupo responde a CQ-E1, CQ-E3 y CQ-G2.

\medskip\noindent\emph{Autorización y derivación.}
\code{ExecutionAuthorization} y \code{DerivedArtifact} cubren la cadena autorización-salida. ODRL aporta acuerdos vinculantes, y PROV-O relaciones generales de generación y derivación, pero ninguno registra el acto específico de autorizar una ejecución concreta de aprendizaje automático ni el artefacto gobernable que esa ejecución produce. La autorización se modela como el acuerdo negociado que vincula permisos con ejecuciones; el artefacto derivado se modela como la salida gobernada de tal ejecución. La invariante INV-1 fuerza la consistencia de esta cadena. El grupo responde a CQ-E2, CQ-E5, CQ-G3 y CQ-G4.

\medskip\noindent\emph{Auditoría y agregación de procedencia.}
\code{AuditEvidence} y \code{CrossParticipantProvenanceRecord} cubren la auditabilidad a través de fronteras administrativas. SPDX proporciona checksums, y PROV-O bundles de procedencia, pero ninguno atiende por sí solo la evidencia orientada a cumplimiento que cruza contextos de participante. La evidencia de auditoría combina un valor de integridad, un firmante y una marca temporal; el registro de procedencia cruzada agrega actividades y entidades relevantes entre contextos de participante, produciendo una narrativa de auditoría que un actor de gobernanza puede recorrer extremo a extremo. El grupo responde a CQ-E4 y CQ-G4.

\medskip\noindent\emph{Contexto de evaluación.}
\code{SharedEvaluationContext} cubre el hueco de comparabilidad que las evaluaciones MLDCAT-AP dejan abierto. MLDCAT-AP registra una medición asociada a un modelo; no reifica la combinación de tarea, versión de dataset, protocolo de evaluación y semilla aleatoria que hace comparables dos mediciones. Las métricas se modelan como recursos vinculados a una evaluación y a un contexto, no como atributos fijos del modelo: la misma versión de modelo puede exhibir precisión, sesgo, latencia o coste computacional distintos bajo datasets, protocolos o infraestructuras distintas. Dos evaluaciones que comparten este contexto son comparables por construcción; evaluaciones que difieren en cualquiera de esas coordenadas no lo son. Esta comparabilidad es operacional y mínima: presupone que la métrica referenciada tiene la misma definición, unidad y procedimiento de cálculo, aspectos que pueden declararse en perfiles de métricas o protocolos cuando el espacio de datos los necesite. La clase responde a CQ-V1, CQ-V2 y CQ-V3.

\paragraph{Defensa de los alineamientos de mayor carga.}\label{sec:alignment-defense}

De los ocho alineamientos \code{rdfs:subClassOf} hacia vocabularios externos, dos llevan significativamente más carga semántica que el resto y merecen defensa explícita porque las contra-lecturas son plausibles y un lector familiarizado con DCAT y ODRL las planteará inmediatamente.

\paragraph{\code{AIAssetOffering} $\sqsubseteq$ \code{dcat:CatalogRecord}, no \code{dcat:Dataset}.} DCAT distingue con cuidado entre el recurso ofrecido (el \code{dcat:Dataset}, que es la cosa sobre la que hay interés) y el registro que describe dicho recurso en un catálogo concreto (el \code{dcat:CatalogRecord}, que es el artefacto de metadatos). DAIMO alinea \code{AIAssetOffering} a \code{dcat:CatalogRecord} porque la oferta es conceptualmente el registro de catalogación del modelo en un espacio de datos concreto, con una política y un conjunto de metadatos específicos de ese contexto, no el modelo en sí ---que se reutiliza directamente como \code{it6:MachineLearningModel}. La contra-lectura que subsumiría \code{AIAssetOffering} a \code{dcat:Dataset} conflaría dos cosas distintas: el modelo (que puede aparecer en varios catálogos con condiciones distintas) y su oferta en un catálogo específico (que es una y con una política propia). La distinción es operacionalmente relevante: una misma \code{it6:MachineLearningModel} puede ser referenciada por varias \code{AIAssetOffering}s con políticas distintas, y la propiedad \code{daimo:offersModel}, alineada a \code{foaf:primaryTopic}, captura que la oferta trata sobre el modelo sin identificarse con él.

\paragraph{\code{ExecutionAuthorization} $\sqsubseteq$ \code{odrl:Agreement}, no \code{prov:wasDerivedFrom odrl:Agreement}.} La autorización de ejecución es el resultado específico de una negociación DSP que vincula permisos ODRL concretos a invocaciones concretas, y ODRL modela precisamente este tipo de instrumento vinculante como \code{odrl:Agreement} ---el equivalente a un contrato suscrito entre un asignador y un asignatario, con permisos y restricciones efectivas. La contra-lectura que relacionaría la autorización con el acuerdo mediante \code{prov:wasDerivedFrom} trataría la autorización como un artefacto distinto del acuerdo, introduciendo una distinción que ODRL no necesita: un \code{Agreement} ya es la entidad vinculante. La subsunción aprovecha toda la semántica ODRL sobre permisos, prohibiciones y obligaciones sin redefinirla, y permite que herramientas ODRL existentes traten las autorizaciones DAIMO como acuerdos estándar. Lo que DAIMO añade sobre \code{odrl:Agreement} es la propiedad asimétrica \code{daimo:authorizesRun}, que vincula el acuerdo a las ejecuciones \code{it6:Run} que cubre: una relación que ODRL no prescribe pero que el puente al plano PROV requiere.

Los cinco alineamientos restantes se justifican a continuación, ordenados desde el más discutible hasta el menos problemático.

\paragraph{\code{ParticipantRole} $\sqsubseteq$ \code{prov:Role}.} PROV-O acota \code{Role} a una asociación cualificada dentro de una actividad, mientras que el rol DAIMO se acota a un contexto de dataspace y normalmente sobrevive a una ejecución concreta. La subsunción es defendible sólo en sentido \emph{is-a-kind-of}: un rol DAIMO es un rol PROV cuando cualifica una actividad específica, y la diferencia de alcance se captura mediante la propiedad que lo vincula a un \code{edc:ParticipantContext}. La alternativa menos comprometida sería \code{skos:related} con \code{prov:Role}; se descarta porque impediría que los razonadores clasifiquen roles DAIMO como roles PROV cuando efectivamente cualifican una asociación. Se acepta la subsunción con la condición documentada de que el contexto DAIMO añade alcance administrativo, no contradicción semántica.

\paragraph{\code{ModelDeployment} $\sqsubseteq$ \code{prov:Entity}.} Un despliegue es un artefacto persistente sobre el que ocurren actividades (carga del modelo, inferencia, logging), no una actividad en sí. Modelarlo como \code{prov:Activity} forzaría marcas \code{startedAtTime}/\code{endedAtTime}, adecuadas para una ejecución pero no para una instancia alojada de larga duración. Alinear el despliegue con \code{prov:Entity} preserva la intuición de que poner un modelo en una infraestructura produce una entidad gobernable sobre la que luego se ejecutan actividades.

\paragraph{\code{DerivedArtifact} $\sqsubseteq$ \code{prov:Entity}, \code{dcat:Resource}.} La doble subsunción refleja dos roles complementarios: el artefacto derivado es una entidad de procedencia producida por una ejecución y, al mismo tiempo, un recurso catalogable que un actor de gobernanza puede describir, listar o distribuir. \code{prov:Entity} y \code{dcat:Resource} son compatibles; DCAT no introduce compromiso alguno que contradiga PROV. El cierre OWL-RL del demostrador confirma que esta doble herencia no produce tipados no intencionados: la verificación de implicación inspecciona todas las superclases inferidas de \code{DerivedArtifact} y no emite advertencias.

\paragraph{\code{CrossParticipantProvenanceRecord} $\sqsubseteq$ \code{prov:Bundle}.} PROV-O ofrece dos constructos próximos: \code{prov:Bundle}, un conjunto nombrado de descripciones de procedencia con atribución propia, y \code{prov:Collection}, un conjunto de entidades sin metasemántica de procedencia. El registro DAIMO corresponde al primero: agrupa descripciones de procedencia a través de varios \code{edc:ParticipantContext}s y puede tener su propia atribución, fecha y evidencia. Alinear con \code{prov:Collection} perdería esa capa meta-provenance; alinear con ambos introduciría redundancia.

\paragraph{\code{AuditEvidence} $\sqsubseteq$ \code{prov:Entity}.} La evidencia es un artefacto producido (hash, firmante y marca temporal), no el acto de producirlo. La actividad de auditar se modela separadamente como una \code{prov:Activity} asociada al actor de cumplimiento; \code{AuditEvidence} es lo que esa actividad genera. Alinear la clase con \code{prov:Activity} conflaría actividad y resultado, y debilitaría la distinción entidad-actividad sobre la que PROV se apoya.

\subsection{Fundamentos axiomáticos}\label{sec:axiomatic}

Una taxonomía por sí sola no fija el significado: sólo nombra los tipos. Detrás de la taxonomía de DAIMO hay una infraestructura axiomática que fija qué es cada clase, qué no es y cómo deben comportarse sus instancias. Tres estratos complementarios componen esta infraestructura: axiomas OWL sobre las clases y propiedades nativas, restricciones SHACL nodales sobre las instancias e invariantes SHACL-SPARQL entre clases. Los párrafos siguientes describen el primer estrato; los estratos SHACL aparecen en §\ref{sec:eval-invariants} como parte de la validación.

En el estrato OWL la disyunción entre los nueve tipos nativos de primer nivel se declara mediante el axioma nombrado \code{daimo:TopLevelKindsDisjointness}, instancia de \code{owl:AllDisjointClasses}. Convertir la disyunción en una entidad nombrada, en lugar de emitir pares anónimos \code{owl:disjointWith}, persigue tres beneficios. Primero, la disyunción puede referenciarse en documentación y reportes de validación. Segundo, un razonador puede citarla como causa específica cuando detecta una contradicción, facilitando el diagnóstico. Tercero, la disyunción entre las cinco subclases de \code{ParticipantRole} queda deliberadamente fuera del axioma: un mismo agente jurídico puede asumir varios roles simultáneamente (por ejemplo, proveedor y operador) sin generar inconsistencia lógica, reflejando la realidad de que los roles son propiedades contextuales del agente y no clases mutuamente exclusivas.

Al núcleo de disyunciones se añaden caracterizaciones sobre las propiedades DAIMO. Veintisiete propiedades se declaran funcionales (\code{owl:FunctionalProperty}) sobre aquellas cuyo dominio es una clase con \emph{shape} nodal, de modo que la cardinalidad máxima queda axiomatizada simultáneamente en la capa lógica y en la capa SHACL correspondiente. Cinco propiedades clave (\code{offersModel}, \code{deploysModel}, \code{authorizesRun}, \code{derivedFromRun} y \code{evidenceOf}) se marcan como asimétricas, capturando formalmente que un despliegue no es el modelo que despliega, que una autorización no coincide con la ejecución que autoriza, y así sucesivamente. Cinco pares inversos explícitos (\code{authorizedBy}~$\leftrightarrow$~\code{authorizesRun}, \code{hasDeployment}~$\leftrightarrow$~\code{deploysModel}, \code{hasDerivedArtifact}~$\leftrightarrow$~\code{derivedFromRun}, \code{hasAuditEvidence}~$\leftrightarrow$~\code{evidenceOf}, \code{hasOffering}~$\leftrightarrow$~\code{offersModel}) habilitan consultas inversas nativas sin obligar al publicador del grafo a materializar ambos sentidos.

El estrato SHACL se divide a su vez en dos familias. En primer lugar, cada clase DAIMO-nativa lleva asociada una \emph{shape} nodal de completitud que prescribe qué propiedades son obligatorias en una instancia bien formada; estas nueve \emph{shapes} son la garantía sintáctica de que un publicador ha suministrado todo lo necesario para que las invariantes semánticas operen sobre el grafo. En segundo lugar, tres \emph{shapes} de conformidad (\code{OfferInDAIMOShape}, \code{MachineLearningModelInDAIMOShape} y \code{RunInDAIMOShape}) establecen qué propiedades de clases reutilizadas son obligatorias cuando un grafo declara adherirse al perfil DAIMO. La distinción entre ambas familias es deliberada: DAIMO no sobrescribe la ontología subyacente, porque eso fragmentaría el ecosistema, sino que declara su compromiso con un subconjunto concreto de propiedades sobre las clases externas, haciendo explícito el criterio de adhesión sin modificar las definiciones originales.

Por encima de las \emph{shapes} se sitúan las seis invariantes SHACL-SPARQL entre clases (INV-1 a INV-6) que se examinan en detalle en §\ref{sec:eval-invariants}. Son el estrato que distingue a DAIMO de un simple perfil descriptivo: codifican reglas cuyo sentido no puede expresarse como restricción local a una sola \emph{shape} porque involucran la coherencia entre varias clases simultáneamente. El Listado~\ref{lst:axioms} muestra cómo se codifican los axiomas OWL para el caso concreto de \code{ExecutionAuthorization}, combinando la subclase con \code{odrl:Agreement}, la asimetría de \code{authorizesRun}, la funcionalidad de su inversa y la pertenencia al axioma de disyunción nombrado.

\begin{lstlisting}[label=lst:axioms,caption={Axiomas centrales alrededor de \code{ExecutionAuthorization}.}]
daimo:ExecutionAuthorization a owl:Class ;
    rdfs:subClassOf odrl:Agreement ;
    skos:definition "An ODRL agreement produced by a DSP contract
        negotiation that binds specific permissions to run invocations..."@en .

daimo:authorizesRun a owl:ObjectProperty , owl:AsymmetricProperty ;
    rdfs:domain daimo:ExecutionAuthorization ;
    rdfs:range  it6:Run .

daimo:authorizedBy a owl:ObjectProperty , owl:FunctionalProperty ;
    owl:inverseOf daimo:authorizesRun ;
    rdfs:domain it6:Run ;
    rdfs:range  daimo:ExecutionAuthorization .

daimo:TopLevelKindsDisjointness a owl:AllDisjointClasses ;
    owl:members ( daimo:AIAssetOffering
                  daimo:ExecutionAuthorization
                  daimo:ModelDeployment
                  daimo:DerivedArtifact
                  daimo:IOContract
                  daimo:AuditEvidence
                  daimo:CrossParticipantProvenanceRecord
                  daimo:SharedEvaluationContext
                  daimo:ParticipantRole ) .
\end{lstlisting}

\paragraph{Superficie de reutilización.}\label{sec:reuse}

La Tabla~\ref{tab:reuse} resume los vocabularios reutilizados y el uso concreto de cada uno. Tres observaciones son importantes. La primera es que las alineaciones hacia DCAT y ODRL son con los vocabularios estándar, no con entidades internas de ninguna implementación: este punto es relevante porque en la práctica de Eclipse EDC las dos capas se confunden con frecuencia. La segunda es que DAIMO mantiene hacia DSP una relación de mapeo informativo mediante \code{skos:closeMatch} y \code{skos:related}, no de subsunción: DSP es un protocolo de mensajería, no una ontología de activos, y subsumir clases DAIMO bajo constructos DSP introduciría tipados espurios. La tercera es que de toda la extensión EDC sólo \code{edc:ParticipantContext} aparece explícitamente en DAIMO, referenciado por su utilidad para anclar la procedencia agregada de un participante.

\begin{table*}[t]
\small\sf\centering
\caption{Vocabularios reutilizados por DAIMO y uso concreto de cada uno.\label{tab:reuse}}
\begin{tabular}{@{}p{2.6cm}p{5cm}p{7cm}@{}}
\toprule
\textbf{Vocabulario} & \textbf{Términos reutilizados} & \textbf{Uso en DAIMO} \\
\midrule
MLDCAT-AP 3.0.0 (\code{it6:}) & \code{MachineLearningModel}, \code{Task}, \code{Run}, \code{Flow}, \code{Evaluation}, \code{EvaluationMeasure}, \code{Benchmark}, \code{ComputerInfrastructure}, \code{Hardware}, \code{Library}, \code{HarmRisk}, \code{Modality}, \code{trainedOn}, \code{testedOn} & Capa del activo de IA; DAIMO no redefine el modelo. \\
DCAT 2 (\code{dcat:}) & \code{Catalog}, \code{Dataset}, \code{Distribution}, \code{DataService}, \code{CatalogRecord}, \code{endpointURL} & Publicación y descubrimiento. \\
ODRL 2.2 (\code{odrl:}) & \code{Offer}, \code{Agreement}, \code{Permission}, \code{Prohibition}, \code{hasPolicy}, \code{target}, \code{assigner}, \code{assignee} & Política, oferta y acuerdo. \\
PROV-O (\code{prov:}) & \code{Activity}, \code{Entity}, \code{Agent}, \code{Bundle}, \code{Role}, \code{wasGeneratedBy}, \code{wasAssociatedWith} & Procedencia y roles. \\
SPDX (\code{spdx:}) & \code{Checksum}, \code{algorithm}, \code{checksumValue} & Integridad de evidencia de auditoría. \\
DSP (\code{dspace:}) & \code{ContractOffer}, \code{ContractNegotiation}, \code{TransferProcess} & Mapeos informativos \code{skos:closeMatch} y \code{skos:related}; no subsunción. \\
EDC (\code{edc:}) & \code{ParticipantContext} & Única extensión específica de EDC referenciada. \\
\bottomrule
\end{tabular}
\end{table*}

Tres subpropiedades que podrían parecer naturales no se declaran, y la ausencia merece documentación explícita. \code{daimo:authorizesRun} no se declara subpropiedad de \code{prov:used}; \code{daimo:grantedTo} no se declara subpropiedad de \code{prov:qualifiedAssociation}; \code{daimo:evidenceOf} no se declara subpropiedad de \code{prov:hadActivity}. En los tres casos, el alineamiento aparentemente intuitivo produciría, bajo razonamiento RDFS u OWL-RL, que las autorizaciones se tipasen como actividades PROV, los agentes receptores como asociaciones reificadas y las evidencias como objetos de influencia. Estos tipados contradicen el modelo conceptual del puente DAIMO. La verificación de implicación (§\ref{sec:eval-engineering}) se diseñó precisamente para detectar errores de este tipo antes de que contaminen el cierre inferido.

% =====================================================================
\section{Evaluación}\label{sec:eval}

La evaluación distingue dos objetos: DAIMO como artefacto ontológico y la evaluación comparativa de modelos que DAIMO permite describir. El primer objeto se examina como evaluación técnica del artefacto: adecuación a requisitos, consistencia formal, conformidad SHACL, ausencia de pitfalls críticos, comportamiento de los alineamientos y respuesta a preguntas de competencia. El segundo aparece sólo dentro del demostrador como comparación de resultados bajo un contexto compartido, sin constituir un benchmark empírico de modelos reales. En términos de LOT, la validación de propósito se cubre de forma acotada mediante un grafo de instancia realista, consultas SPARQL ejecutables y pruebas negativas. Siguiendo la distinción habitual entre cobertura de requisitos, corrección formal, validación basada en restricciones y utilidad aplicada~\cite{brank05}, preguntamos si el perfil representa el intercambio objetivo, si el artefacto OWL es formalmente coherente, si las restricciones SHACL discriminan grafos válidos e inválidos, y si las preguntas de competencia son ejecutables sobre el grafo materializado. Sobre esta distinción, el análisis sigue el patrón de cinco dimensiones de RePlanIT~\cite{kurteva}: instanciación del escenario (§\ref{sec:eval-usecase}), verificación formal mediante razonamiento, comprobación de implicación y escaneo OOPS! (§\ref{sec:eval-engineering}), validación de instancias frente a SHACL (§\ref{sec:eval-invariants}), ejecución de las veintitrés preguntas de competencia (§\ref{sec:eval-cqs}) y verificación de los alineamientos de reutilización (§\ref{sec:eval-reuse}). Las cinco dimensiones son reproducibles con el paquete de evaluación que acompaña al artefacto.

\subsection{Modelado del caso de uso}\label{sec:eval-usecase}

El escenario introducido en §\ref{sec:scenario} se instancia sobre DAIMO como demostrador controlado. El grafo resultante contiene 225 tripletas de datos que ejercen al menos una vez cada clase nativa de la ontología y sobre las que se ejecutan las veintitrés consultas de competencia reportadas en §\ref{sec:eval-cqs}. Todas las organizaciones, nombres de modelo, versiones de dataset, métricas y endpoints son anónimos o sintéticos. El demostrador no es un estudio de despliegue, un \emph{gold standard} ni una validación con catálogos reales de modelos; su función es probar adecuación representacional y hacer ejecutables los artefactos de validación.

El Proveedor~A publica tres ofertas gobernadas en un catálogo federado: dos versiones de un mismo modelo geoespacial de predicción de riesgo y una variante de sólo lectura con política más estricta. Las ofertas comparten tarea y dominio, pero difieren en las condiciones ODRL asociadas. Esta parte del escenario ejercita la distinción entre el modelo como activo de IA descrito con MLDCAT-AP y la oferta como registro contextual de catálogo con condiciones de uso propias.

El modelo seleccionado se despliega sobre un entorno de cómputo descrito y se expone mediante dos interfaces de servicio: un endpoint REST con entrada JSON y autenticación por bearer token, y un endpoint gRPC con protobuf y autenticación mTLS. El uso multi-endpoint ejercita la decisión de modelado según la cual un mismo despliegue puede exponer varios servicios y contratos, en lugar de forzar un despliegue distinto por endpoint.

El Consumidor municipal~B obtiene una autorización de ejecución con fechas de vigencia explícitas, invoca el modelo y recibe un artefacto de predicción derivado. La ejecución se vincula con evidencia de auditoría que contiene registro de integridad, firmante y marca temporal, de modo que un actor de gobernanza pueda reconstruir cómo se produjo el artefacto y bajo qué autorización. Un registro de procedencia cruzada agrega la cadena a través de los contextos de proveedor, consumidor y plataforma. El Evaluador~D publica dos resultados sintéticos de evaluación bajo el mismo benchmark, protocolo y semilla aleatoria, permitiendo comparar dos versiones del modelo bajo condiciones compartidas. La Figura~\ref{fig:scenario} ilustra el escenario instanciado.

\begin{figure*}[t]
\centering
\begin{tikzpicture}[node distance=0.6cm and 0.8cm,
    instance/.style={rectangle, draw=black!60, fill=yellow!15, rounded corners=2pt,
      minimum width=2.2cm, minimum height=0.6cm, align=center, font=\ttfamily\scriptsize,
      inner sep=2pt},
    extinstance/.style={rectangle, draw=black!60, fill=blue!8, rounded corners=2pt,
      minimum width=2.2cm, minimum height=0.6cm, align=center, font=\ttfamily\scriptsize,
      inner sep=2pt},
    actorbox/.style={rectangle, draw=black!70, fill=gray!15, thick, rounded corners=3pt,
      minimum width=1.6cm, minimum height=0.55cm, align=center, font=\sffamily\scriptsize\bfseries,
      inner sep=3pt}]

  \node[actorbox] (upm) {Proveedor~A};
  \node[actorbox, right=1.0cm of upm] (leg) {Consumidor~B};
  \node[actorbox, right=1.0cm of leg] (inesdata) {Plataforma~C};
  \node[actorbox, right=1.0cm of inesdata] (csic) {Evaluador~D};
  \node[actorbox, right=1.0cm of csic] (gaia) {Cumplimiento~E};

  \node[extinstance, below=0.8cm of upm] (catalog) {cat\'alogo proveedor\\\scriptsize (dcat:Catalog)};
  \node[instance, below=0.6cm of catalog, xshift=-1.8cm] (off1) {oferta v2};
  \node[instance, below=0.6cm of catalog] (off2) {oferta v1};
  \node[instance, below=0.6cm of catalog, xshift=1.8cm] (off3) {oferta s\'olo lectura};

  \node[extinstance, below=0.8cm of off2] (model) {modelo riesgo v2\\\scriptsize (MLModel)};
  \node[instance, below=0.6cm of model] (deploy) {despliegue modelo};
  \node[extinstance, below=0.5cm of deploy, xshift=-1.6cm] (rest) {REST svc};
  \node[extinstance, below=0.5cm of deploy, xshift=1.6cm] (grpc) {gRPC svc};

  \node[instance, below=1.5cm of leg, xshift=4cm] (auth) {autorizaci\'on};
  \node[extinstance, below=0.6cm of auth] (run) {ejecuci\'on consumidor\\\scriptsize (it6:Run)};
  \node[instance, below=0.6cm of run, xshift=-1.6cm] (derived) {predicci\'on derivada};
  \node[instance, below=0.6cm of run, xshift=1.6cm] (audit) {evidencia auditor\'ia};

  \node[instance, below=0.8cm of derived, xshift=1.6cm] (bundle) {bundle cruzado};

  \node[extinstance, below=1.2cm of csic] (eval2) {evaluaci\'on v2 (0.89)};
  \node[extinstance, below=0.3cm of eval2] (eval1) {evaluaci\'on v1 (0.82)};
  \node[instance, below=0.3cm of eval1] (ctx) {contexto compartido};

  \draw[objpropertyarrow] (upm) -- node[left, font=\tiny\sffamily] {publishes} (catalog);
  \draw[objpropertyarrow] (catalog) -- (off1);
  \draw[objpropertyarrow] (catalog) -- (off2);
  \draw[objpropertyarrow] (catalog) -- (off3);
  \draw[objpropertyarrow] (off2) -- node[right, font=\tiny\sffamily] {offersModel} (model);
  \draw[objpropertyarrow] (off1) to[bend right=15] (model.north west);
  \draw[objpropertyarrow] (off3) to[bend left=15] (model.north east);

  \draw[objpropertyarrow] (model) -- node[right, font=\tiny\sffamily] {hasDeployment} (deploy);
  \draw[objpropertyarrow] (deploy) -- (rest);
  \draw[objpropertyarrow] (deploy) -- (grpc);

  \draw[objpropertyarrow] (leg) -- node[right, font=\tiny\sffamily] {obtains} (auth);
  \draw[objpropertyarrow] (auth) -- node[right, font=\tiny\sffamily] {authorizesRun} (run);
  \draw[objpropertyarrow] (run) -- (derived);
  \draw[objpropertyarrow] (run) -- (audit);
  \draw[objpropertyarrow] (run) to[bend right=20] node[above, font=\tiny\sffamily] {onDeployment} (deploy.east);

  \draw[objpropertyarrow] (derived) -- (bundle);
  \draw[objpropertyarrow] (audit) -- (bundle);

  \draw[objpropertyarrow] (csic) -- node[right, font=\tiny\sffamily] {publishes} (eval2);
  \draw[objpropertyarrow] (eval2) -- (ctx);
  \draw[objpropertyarrow] (eval1) -- (ctx);
  \draw[objpropertyarrow, dashed] (model) -- node[above, font=\tiny\sffamily, sloped] {evaluatedBy} (eval2);

  \draw[objpropertyarrow, dashed] (gaia) -- node[right, font=\tiny\sffamily] {audits} (bundle.north east);
  \draw[objpropertyarrow, dashed] (inesdata) -- node[left, font=\tiny\sffamily] {hosts} (deploy.north east);

\end{tikzpicture}
\caption{Demostrador anónimo instanciado sobre DAIMO. Los rect\'angulos grises redondeados de la fila superior son los cinco contextos de participante; las instancias sombreadas en amarillo son constructos de gobernanza introducidos por DAIMO; las instancias en azul claro son recursos descritos con vocabularios reutilizados. Las flechas s\'olidas denotan relaciones de dominio; las discontinuas muestran el rol funcional del participante sobre la instancia inferior.}
\label{fig:scenario}
\end{figure*}

\subsection{Verificación formal}\label{sec:eval-engineering}

El razonador OWL~2~DL HermiT, invocado mediante \code{owlready2}, reporta que la unión del núcleo, los alineamientos y el grafo de ejemplo es lógicamente consistente y que no existen clases insatisfactibles; el tiempo de razonamiento reportado es de 1,25 segundos. La materialización OWL-RL (biblioteca \code{owlrl}) deriva 1\,170 tripletas adicionales sobre 818 iniciales (1\,988 totales) en 0,58 segundos, sin producir individuos tipados como \code{owl:Nothing} ni subclases insatisfactibles. Ambos razonadores coinciden, con HermiT como validador de la corrección lógica y OWL-RL como materializador pragmático para las consultas que requieren inferencia por subpropiedad o subclase.

La consistencia lógica es condición necesaria pero no suficiente para la corrección semántica de una ontología de integración. Una axiomatización puede ser perfectamente consistente y, al mismo tiempo, implicar tipados no intencionados si no declara disyunción explícita con la clase adecuada o si adopta una subpropiedad cuyo rango fuerza una subsunción indeseada. Por eso DAIMO implementa adicionalmente una \emph{verificación de implicación} sobre el cierre OWL-RL: para cada clase nativa, la comprobación enumera todas las superclases inferidas y las contrasta con una lista de superclases prohibidas. Por ejemplo, \code{AuditEvidence} no debe ser inferida como \code{prov:Activity}, porque es una cosa producida por actividades, no una actividad. El reporte cubre las catorce clases nativas (nueve de primer nivel más cinco subclases de \code{ParticipantRole}) y no emite ninguna advertencia. Esta comprobación es complementaria a la de consistencia y captura errores que un razonador DL acepta porque son lógicamente consistentes aunque semánticamente incorrectos.

El escáner OOPS!~\cite{oops14}, consultado por API REST sobre la unión del núcleo y los alineamientos, reporta cero pitfalls Críticos, cero Importantes y dos Menores. El Menor P13 (``relaciones inversas no declaradas'') afecta a 34 elementos, y corresponde a propiedades para las que no existe una inversa semánticamente significativa (por ejemplo, \code{daimo:integrityHash}: nombrar una inversa ``tiene hash'' no añade capacidad de consulta). El Menor P04 (``elementos no conectados'') afecta a siete elementos que son clases externas declaradas localmente para permitir que el razonador las interprete sin exigir que el archivo las defina completamente. Ambos Menores son características habituales de perfiles de integración y no indican defectos de modelado; su aparición es precisamente el resultado de la disciplina \emph{reuse-first} que el perfil aplica.

Las métricas de calidad estructural se interpretan como señales diagnósticas, no como criterios absolutos de aceptación. Conteos como profundidad de jerarquía, amplitud taxonómica, conectividad, equivalencias explícitas o elementos aislados ayudan a localizar zonas que merecen inspección, pero no dictan por sí mismos una corrección automática. En OWL, un ciclo de \code{rdfs:subClassOf} puede expresar equivalencia implícita antes que un error; una \code{owl:equivalentClass} puede ser un alineamiento legítimo; y una disyunción sólo debe añadirse cuando el compromiso ontológico es verdadero en todo modelo posible del dominio. DAIMO usa por ello OOPS!, HermiT, OWL-RL, SHACL, pruebas negativas y preguntas de competencia como capas complementarias: cada una observa un tipo de riesgo distinto y ninguna sustituye el juicio de modelado.

\subsection{Validación SHACL: estructural, invariantes y banco negativo}\label{sec:eval-invariants}

DAIMO superpone tres familias de \emph{shapes} SHACL, cada una respondiendo a una pregunta de validación distinta. Las \emph{shapes} estructurales preguntan si toda instancia DAIMO-nativa porta las propiedades obligatorias de su clase. Las \emph{shapes} de conformidad preguntan si las clases reutilizadas externas satisfacen las obligaciones adicionales que el perfil les impone. Las invariantes SPARQL preguntan si las relaciones entre oferta, política, despliegue, autorización, ejecución, artefacto derivado y evidencia siguen siendo coherentes. Juntas, las tres familias llevan la validación de lo sintáctico a lo semántico: de ``¿está esta instancia bien formada?'' a ``¿es este grafo internamente consistente?''. La capa SHACL funciona así como contrato de datos del perfil DAIMO; no introduce una ontología paralela ni exige trazabilidad, formatos de modelo o métricas universales que no estén justificadas por las preguntas de competencia.

En términos de implementación, las restricciones se expresan como \code{sh:NodeShape} dirigidas a clases DAIMO o clases reutilizadas bajo el perfil, y como \code{sh:PropertyShape} internas que fijan cardinalidades, tipos de dato, clases esperadas, patrones o enumeraciones cerradas. En un flujo operativo, un proveedor podría publicar metadatos RDF o JSON-LD de un modelo, un conector o servicio de catálogo ejecutaría la capa \code{daimo-shapes.ttl}, y el \code{sh:ValidationReport} resultante indicaría si el grafo conforma o qué propiedad concreta incumple el perfil. Este flujo describe el modo esperado de consumo de las \emph{shapes}; no se presenta como un conector productivo ya desplegado ni como una validación de la ontología TBox en sí misma.

Las nueve \emph{shapes} estructurales imponen completitud mínima por clase; el grafo positivo de 225 tripletas satisface \code{conforms=True}. Las tres \emph{shapes} de conformidad añaden obligaciones operativas sobre clases reutilizadas: \code{OfferInDAIMOShape} exige que toda \code{odrl:Offer} utilizada en DAIMO declare \code{odrl:assigner} y \code{odrl:target} (a nivel de política o por permiso, mediante el conectivo \code{sh:or}); \code{MachineLearningModelInDAIMOShape} exige política ODRL, título e identificador sobre \code{it6:MachineLearningModel}; \code{RunInDAIMOShape} exige \emph{flow}, algoritmo, agente asociado y marca temporal sobre \code{it6:Run}. Adicionalmente, \code{sh:in} restringe los valores de \code{daimo:authMethod} a un vocabulario cerrado de nueve tokens (\code{none}, \code{api-key}, \code{basic}, \code{bearer}, \code{oauth2-bearer}, \code{oauth2-client-credentials}, \code{mtls}, \code{jwt}, \code{dataspace-token}), y \code{sh:pattern} restringe \code{daimo:protocol} a un patrón que acepta los nombres habituales de protocolo de evaluación (\code{holdout}, \code{stratified-holdout}, \code{leave-one-out}, \code{temporal-split}, \code{<n>-fold-cv}, \code{stratified-<n>-fold-cv}, \code{bootstrap-<n>}).

Sobre este estrato estructural se sitúan las seis invariantes SHACL-SPARQL entre clases, codificadas mediante \code{sh:SPARQLConstraint}, que son el elemento diferenciador de DAIMO frente a un perfil puramente descriptivo. La invariante INV-1 garantiza la consistencia de la cadena de autorización de derivación: todo \code{DerivedArtifact} debe portar una propiedad \code{daimo:underAuthorization} cuya autorización cubra precisamente la ejecución de la que el artefacto se deriva. La INV-2 impone que el agente asociado a una ejecución mediante \code{prov:wasAssociatedWith} aparezca también como \code{daimo:grantedTo} en alguna de las autorizaciones que cubren esa ejecución, haciendo explícito que el actor registrado en PROV y el beneficiario de la autorización ODRL deben ser la misma entidad. La INV-3 verifica la coherencia del plano de datos: el \code{dcat:DataService} expuesto por un \code{ModelDeployment} debe servir exactamente el modelo que el despliegue declara desplegar, eliminando la posibilidad de que un endpoint anuncie un modelo pero sirva otro. La INV-4 exige que el \code{daimo:expiresAt} de cada \code{ExecutionAuthorization} sea estrictamente posterior al \code{prov:startedAtTime} de toda ejecución que la autorización cubre, reflejando que una autorización no puede justificar ejecuciones que comenzaron tras su caducidad. La INV-5 requiere que el modelo referenciado por \code{daimo:offersModel} en una \code{AIAssetOffering} aparezca como \code{odrl:target} de la política asociada, a nivel de política o en alguno de sus permisos; si la política no gobierna el modelo registrado, la cadena de gobernanza está rota aunque ambos recursos existan. La INV-6 cierra el mismo triángulo por el lado del proveedor: el \code{daimo:offeredBy} de toda oferta debe coincidir con el \code{odrl:assigner} de su política, impidiendo inconsistencias entre la atribución del registro de catálogo y la autoría declarada en la oferta ODRL. El grafo positivo conforma las seis invariantes sin advertencias. El Listado~\ref{lst:inv5} muestra la codificación concreta de INV-5, ilustrando el patrón \code{FILTER NOT EXISTS} que captura la disyunción entre política global y permiso individual.

\begin{lstlisting}[label=lst:inv5,caption={Invariante INV-5 como \code{sh:SPARQLConstraint}.}]
daimo:OfferingPolicyTargetInvariant a sh:NodeShape ;
    sh:targetClass daimo:AIAssetOffering ;
    sh:sparql [
        a sh:SPARQLConstraint ;
        sh:message "INV-5 violation: offering's offersModel does not
                    appear as odrl:target of the attached policy."@en ;
        sh:prefixes daimo:_invariantPrefixes ;
        sh:select """
            SELECT $this ?model ?policy WHERE {
                $this daimo:offersModel ?model ;
                      daimo:hasOfferPolicy ?policy .
                FILTER NOT EXISTS { ?policy odrl:target ?model }
                FILTER NOT EXISTS { ?policy odrl:permission/odrl:target ?model }
            }
        """ ;
    ] .
\end{lstlisting}

Las \emph{shapes} que conforman sobre el grafo positivo son necesarias pero no suficientes para evaluar si las restricciones son semánticamente discriminantes: una restricción trivialmente satisfacible sobre cualquier grafo pasaría la validación positiva sin aportar garantías reales. Por eso DAIMO incorpora un grafo negativo de 118 tripletas que codifica seis casos deliberadamente inválidos ---\code{bad:INV1-artifact}, \code{bad:INV2-run}, \code{bad:INV3-deployment}, \code{bad:INV4-auth}, \code{bad:INV5-offering} y \code{bad:INV6-offering}---, cada uno focalizado en una invariante. La prueba carga la ontología, los \emph{shapes} y el grafo negativo, ejecuta SHACL y verifica dos condiciones: que la validación falle globalmente y que cada nodo focalizado esperado aparezca en el informe con el mensaje específico de la invariante violada. Las seis invariantes se disparan correctamente sobre su violación correspondiente. Esta evidencia, complementaria a la validación positiva, respalda que las invariantes capturan reglas de negocio reales.

\subsection{Preguntas de competencia}\label{sec:eval-cqs}

Las veintitrés preguntas de competencia se expresan como consultas SPARQL y se ejecutan sobre el cierre OWL-RL materializado del grafo positivo. La Tabla~\ref{tab:cqs} resume la cobertura por categoría, los conteos observados y las consultas que requieren razonamiento explícito; el conjunto completo de consultas aparece en el Apéndice~\ref{app:sparql}. Las veintitrés devuelven al menos una fila, lo que verifica que el grafo de demostración contiene la información necesaria para responder a cada pregunta con las primitivas que DAIMO declara. Este resultado se interpreta como validación funcional de cobertura, no como prueba única de validez ontológica: se complementa con HermiT, OWL-RL, OOPS!, SHACL positivo y pruebas negativas. Del mismo modo, la longitud de una consulta o su número de uniones se trata como señal diagnóstica de inspección, no como evidencia automática de sobre-ingeniería.

\begin{table*}[t]
\small\sf\centering
\caption{Cobertura de las 23 preguntas de competencia. Cada CQ devuelve el número de filas indicado sobre el grafo positivo cerrado con OWL-RL. Las celdas en negrita señalan las consultas que requieren razonamiento explícito más allá de SPARQL básico.\label{tab:cqs}}
\begin{tabular}{@{}p{3.3cm}p{2.6cm}p{0.8cm}p{6.3cm}@{}}
\toprule
\textbf{Categoría} & \textbf{Consultas (filas)} & \textbf{N.} & \textbf{Inferencia requerida} \\
\midrule
Registro y publicación (R) & R1(3) R2(1) R3(1) R4(2) R5(2) & 5 & \textbf{R2} (cadena \code{subPropertyOf}), \textbf{R5} (ruta \code{subClassOf}) \\
Descubrimiento y selección (D) & D1(3) D2(2) D3(4) D4(1) & 4 & \textbf{D1} (subtipo de tarea), \textbf{D2} (coincidencia de política) \\
Ejecución y auditabilidad (E) & E1(4) E2(1) E3(2) E4(1) E5(1) & 5 & \textbf{E1} (\code{offersModel}~$\sqsubseteq$~\code{foaf:primaryTopic}), \textbf{E2} (unión agreement-run) \\
Evaluación y reproducibilidad (V) & V1(1) V2(1) V3(2) V4(1) V5(4) & 5 & \textbf{V2}, \textbf{V3} (agregación) \\
Puente de gobernanza (G) & G1(1) G2(2) G3(1) G4(1) & 4 & \textbf{G4} (agregación \code{GROUP\_CONCAT}); G3 ejercita directamente \code{ExecutionAuthorization} como puente DAIMO--ODRL \\
\bottomrule
\end{tabular}
\end{table*}

\begin{table*}[t]
\scriptsize\sf\centering
\caption{Trazabilidad resumida CQ $\to$ clases $\to$ propiedades $\to$ resultado. La tabla permite inspeccionar dentro del PDF qué fragmento semántico ejerce cada CQ.\label{tab:cq-trace}}
\begin{tabular}{@{}p{0.8cm}p{3.8cm}p{4.0cm}p{4.9cm}r@{}}
\toprule
\textbf{CQ} & \textbf{Pregunta resumida} & \textbf{Clases principales} & \textbf{Propiedades recorridas} & \textbf{Filas} \\
\midrule
R1 & Modelos publicados por un proveedor & \code{AIAssetOffering}, \code{MachineLearningModel} & \code{offeredBy}, \code{offersModel} & 3 \\
R2 & Oferta, proveedor y fecha de un modelo & \code{AIAssetOffering}, \code{CatalogRecord} & \code{offersModel} $\sqsubseteq$ \code{foaf:primaryTopic}, \code{offeredBy}, \code{dct:issued} & 1 \\
R3 & Licencia y política aplicables & \code{MachineLearningModel}, \code{odrl:Offer} & \code{hasOfferPolicy}, \code{dct:license}, \code{odrl:permission} & 1 \\
R4 & Contrato de entrada/salida & \code{ModelDeployment}, \code{IOContract} & \code{hasIOContract}, \code{inputFormat}, \code{outputFormat}, \code{authMethod} & 2 \\
R5 & Rol concreto del proveedor & \code{ParticipantRole}, subclases de rol & \code{offeredBy}, \code{rdf:type/}\newline\code{rdfs:subClassOf*} & 2 \\
D1 & Modelos por tarea y dominio & \code{MachineLearningModel}, \code{Task} & \code{it6:hasTask}, \code{dcat:theme} & 3 \\
D2 & Modelos por patrón de política & \code{AIAssetOffering}, \code{odrl:Offer} & \code{hasOfferPolicy}, \code{odrl:permission}, \code{odrl:prohibition} & 2 \\
D3 & Endpoints y autenticación & \code{ModelDeployment}, \code{DataService}, \code{IOContract} & \code{exposedAs}, \code{endpointURL}, \code{authMethod} & 4 \\
D4 & Umbral mínimo de métrica & \code{Evaluation}, \code{SharedEvaluationContext} & \code{usesEvaluationContext}, \code{it6:hasEvaluationMeasure}, \code{it6:hasValue} & 1 \\
E1 & Endpoint desde una oferta & \code{AIAssetOffering}, \code{ModelDeployment}, \code{IOContract} & \code{offersModel}, \code{hasDeployment}, \code{exposedAs}, \code{hasIOContract} & 4 \\
E2 & Ejecuciones autorizadas & \code{ExecutionAuthorization}, \code{Run} & \code{authorizesRun}, \code{grantedTo}, \code{prov:wasAssociatedWith} & 1 \\
E3 & Implementación e infraestructura & \code{Run}, \code{Flow}, \code{ComputerInfrastructure} & \code{it6:hasFlow}, \code{it6:hasAlgorithm}, \code{onInfrastructure} & 2 \\
E4 & Evidencia de auditoría & \code{AuditEvidence}, \code{spdx:Checksum} & \code{evidenceOf}, \code{integrityHash}, \code{signedBy}, \code{recordedAt} & 1 \\
E5 & Artefactos derivados & \code{DerivedArtifact}, \code{ExecutionAuthorization} & \code{derivedFromRun}, \code{underAuthorization} & 1 \\
V1 & Contexto compartido & \code{SharedEvaluationContext}, \code{Evaluation} & \code{contextTask}, \code{contextDataset}, \code{protocol}, \code{randomSeed} & 1 \\
V2 & Mejor modelo bajo contexto & \code{Evaluation}, \code{MachineLearningModel} & \code{usesEvaluationContext}, \code{it6:hasEvaluationMeasure}, \code{it6:hasValue}, \code{ORDER BY} & 1 \\
V3 & Ranking bajo misma métrica & \code{Evaluation}, \code{SharedEvaluationContext}, \code{MachineLearningModel} & \code{usesEvaluationContext}, \code{it6:evaluates}, \code{it6:hasEvaluationMeasure}, \code{it6:hasValue}, \code{ORDER BY} & 2 \\
V4 & Benchmarks del modelo & \code{MachineLearningModel}, \code{Benchmark} & \code{it6:hasBenchmark}, \code{dct:title} opcional & 1 \\
V5 & Artefactos de reproducibilidad & \code{Flow}, \code{AuditEvidence}, \code{DerivedArtifact} & \code{records}, \code{evidenceOf}, \code{integrityHash} & 4 \\
G1 & Ofertas que incluyen un modelo & \code{AIAssetOffering}, \code{MachineLearningModel} & \code{offersModel}, \code{hasOfferPolicy} & 1 \\
G2 & Despliegues por modelo & \code{ModelDeployment}, \code{IOContract} & \code{deploysModel}, \code{onInfrastructure}, \code{hasIOContract} & 2 \\
G3 & Acuerdo que autorizó una ejecución & \code{ExecutionAuthorization}, \code{Run}, \code{odrl:Agreement} & \code{authorizesRun}, \code{grantedTo}, \code{expiresAt} & 1 \\
G4 & Bundle completo de procedencia & \code{CrossParticipant}\newline\code{ProvenanceRecord}, \code{DerivedArtifact} & \code{records}, \code{spansParticipantContext}, \code{GROUP\_CONCAT} & 1 \\
\bottomrule
\end{tabular}
\end{table*}

Dos resultados conviene señalar. Las consultas CQ-R2 y CQ-E1 dependen de la relación \code{subPropertyOf} entre \code{daimo:offersModel} y \code{foaf:primaryTopic} y no devuelven filas sobre un motor SPARQL puro sin razonamiento. La dependencia es documentada, no un defecto: alinear \code{offersModel} bajo \code{foaf:primaryTopic} hace las ofertas descubribles mediante la propiedad DCAT estándar, beneficio que sólo se materializa bajo inferencia. Un consumidor que necesite ejecutar esas consultas sobre un endpoint puro puede materializar el cierre previamente o reescribir las consultas enumerando ambas propiedades. El segundo resultado es que CQ-G2 devuelve dos filas ---una por endpoint del despliegue multi-protocolo REST y gRPC---, lo que confirma el modelado de múltiples servicios por despliegue como patrón legítimo y directamente consultable desde SPARQL.

Dos fragmentos ilustran el patrón de formalización. El Listado~\ref{lst:cq-g3} muestra CQ-G3: la consulta recupera la autorización que cubre una ejecución concreta, el beneficiario autorizado y la fecha de expiración. La alineación \code{ExecutionAuthorization}~$\sqsubseteq$~\code{odrl:Agreement} hace que ese recurso sea interpretable como acuerdo ODRL sin duplicar propiedades de contrato en la consulta. El Listado~\ref{lst:cq-v3} muestra CQ-V3: dos evaluaciones se comparan sólo porque comparten el mismo \code{SharedEvaluationContext}; el ordenamiento por valor de métrica es significativo dentro de ese contexto y no como ranking global independiente del dataset, protocolo o semilla.

\begin{lstlisting}[label=lst:cq-g3,caption={CQ-G3: autorización de una ejecución concreta.}]
SELECT ?auth ?grantee ?expires WHERE {
  ?auth a daimo:ExecutionAuthorization ;
        daimo:authorizesRun
          <https://example.org/daimo-scenario/run-2026-04-20-legs> ;
        daimo:grantedTo ?grantee ;
        daimo:expiresAt ?expires .
}
\end{lstlisting}

\begin{lstlisting}[label=lst:cq-v3,caption={CQ-V3: ranking de modelos bajo un contexto compartido.}]
SELECT ?model ?value WHERE {
  ?eval a it6:Evaluation ;
        it6:evaluates ?model ;
        daimo:usesEvaluationContext
          <https://example.org/daimo-scenario/evalctx-climatebench-v1-2026-1-holdout> ;
        it6:hasEvaluationMeasure ?m .
  ?m it6:hasValue ?value .
}
ORDER BY DESC(?value)
\end{lstlisting}

\paragraph{Reutilización verificada.}\label{sec:eval-reuse}

La reutilización se verifica directamente sobre el módulo de alineamientos. Ocho \code{rdfs:subClassOf} y cinco \code{rdfs:subPropertyOf}, listados en la Tabla~\ref{tab:classes} y visualizados en la Figura~\ref{fig:classes}, conectan DAIMO con sus cinco vocabularios parentales (DCAT, MLDCAT-AP, ODRL, PROV-O y FOAF). La sexta subpropiedad DAIMO hacia un vocabulario externo, \code{hasOffering}~$\sqsubseteq$~\code{foaf:isPrimaryTopicOf}, se declara en el núcleo porque implementa la navegación inversa de \code{offersModel}. La verificación opera en dos dimensiones.

La primera es \emph{estructural}: toda clase nativa que podría haberse definido aisladamente se sitúa bajo un padre externo siempre que dicho padre exista, y toda propiedad nueva extiende una propiedad reutilizada siempre que ésta ya cargue el significado previsto. No se introduce ninguna clase por razones estilísticas, y las clases \code{IOContract} y \code{SharedEvaluationContext} quedan sin padre externo sólo porque ningún padre con semántica veraz existe (§\ref{sec:classes}).

La segunda es \emph{semántica}: los alineamientos no deben producir inferencias no intencionadas. La verificación de implicación de §\ref{sec:eval-engineering} inspecciona el cierre OWL-RL de cada alineamiento y confirma que ninguno produce una superclase semánticamente incorrecta sobre ninguna clase nativa. El razonador reporta cero advertencias sobre las catorce clases. Tres alineamientos de subpropiedad hacia PROV-O se excluyeron deliberadamente del módulo por el mismo motivo: cada uno habría producido un tipado prohibido bajo cierre RDFS (§\ref{sec:reuse}).

\subsection{Cobertura de requisitos no funcionales}\label{sec:eval-nfr}

Los requisitos no funcionales se evalúan sobre el artefacto ontológico y sus flujos de validación, no sobre una plataforma de ejecución federada. En particular, la escalabilidad se interpreta como crecimiento del número de instancias de intercambio gobernado en el ABox; el núcleo DAIMO permanece estable en catorce clases nativas. Por tanto, el resultado no mide concurrencia, latencia de un conector ni rendimiento productivo de inferencia en línea, sino tratabilidad reproducible de grafos DAIMO crecientes bajo las operaciones que el perfil exige: parseo RDF, materialización OWL-RL, validación SHACL y consultas SPARQL.

\begin{table*}[t]
\scriptsize\sf\centering
\caption{Cobertura de requisitos no funcionales de DAIMO. La columna de estado distingue evidencia reproducible del artefacto y riesgo residual.\label{tab:nfr-coverage}}
\begin{tabular}{@{}p{2.7cm}p{8.0cm}p{4.0cm}@{}}
\toprule
\textbf{Requisito} & \textbf{Evidencia reproducible} & \textbf{Estado y riesgo residual} \\
\midrule
Consistencia lógica & HermiT reporta consistencia y cero clases insatisfactibles; OWL-RL materializa 1\,170 tripletas sin \code{owl:Nothing}; la verificación de implicaciones inspecciona catorce clases nativas con cero advertencias. & Cumplimiento fuerte. La consistencia se complementa con pruebas negativas SHACL para evitar restricciones triviales. \\
Interoperabilidad & El perfil usa RDF/OWL~2~DL, SHACL y SPARQL; reutiliza DCAT, MLDCAT-AP, ODRL, PROV-O, FOAF, SKOS y SPDX mediante IRIs externos y alineamientos axiomáticos. & Cobertura fuerte en el artefacto. La interoperabilidad operacional dependerá de la publicación y consumo efectivos en infraestructuras de terceros. \\
Escalabilidad & El benchmark genera grafos sintéticos conformes de 100 y 1\,000 unidades de intercambio; ambos conforman SHACL y responden las consultas de control. & Cobertura medida y acotada. SHACL es el coste dominante; no se afirma rendimiento productivo ni resultado de 10\,000 unidades sin ejecutarlo. \\
Extensibilidad & El núcleo mantiene alcance mínimo y versionado SemVer; las extensiones se orientan a perfiles complementarios, por ejemplo \code{daimo-discovery}, sin absorber lógica de ranking o ejecución. & Cumplimiento fuerte para evolución controlada. Nuevos conceptos deben justificar que no existen en vocabularios reutilizados. \\
Modularidad & Núcleo, alineamientos, \emph{shapes}, ejemplos, consultas, pruebas negativas y reportes están separados; cada módulo puede ejecutarse o revisarse de forma independiente. & Cobertura fuerte. La independencia modular reduce el riesgo de romper el núcleo al añadir perfiles externos. \\
Documentación & Cada término incluye etiquetas y definiciones SKOS; la documentación suplementaria traza clases, propiedades, CQs, invariantes y comandos de reproducción. & Cobertura fuerte para revisión local y reproducción. La calidad editorial de la documentación pública deberá mantenerse sincronizada con futuras versiones del perfil. \\
\bottomrule
\end{tabular}
\end{table*}

\begin{table*}[t]
\small\sf\centering
\caption{Benchmark sintético de escalabilidad acotada. Cada unidad contiene una oferta, modelo, política ODRL, despliegue, servicio, contrato de I/O, autorización, ejecución, artefacto derivado, evidencia de auditoría, evaluación y bundle de procedencia.\label{tab:scalability}}
\begin{tabular}{@{}r r r r r r r c@{}}
\toprule
\textbf{Unidades} & \textbf{Tripletas datos} & \textbf{Tripletas cierre} & \textbf{Parseo} & \textbf{OWL-RL} & \textbf{SHACL} & \textbf{SPARQL} & \textbf{Conforma} \\
 & & & \textbf{(s)} & \textbf{(s)} & \textbf{(s)} & \textbf{(s)} & \\
\midrule
100 & 8\,053 & 16\,248 & 0,468 & 5,478 & 10,374 & 0,080 & Sí \\
1\,000 & 80\,053 & 147\,648 & 4,660 & 53,672 & 135,010 & 0,356 & Sí \\
\bottomrule
\end{tabular}
\end{table*}

% =====================================================================
\section{Discusión}\label{sec:discussion}

DAIMO articula publicación, política, ejecución, procedencia y evaluación en un perfil de integración para el intercambio gobernado de modelos a través de fronteras organizativas. Su contribución es semántica y de gobernanza: no estandariza pesos, arquitecturas internas ni mecanismos de inferencia, sino que hace consultable y validable la relación entre metadatos de modelo, políticas, despliegues, ejecuciones, evidencias y contextos de evaluación. Esta posición la diferencia tanto de ontologías de dominio publicadas recientemente ---RePlanIT~\cite{kurteva}, GloSIS~\cite{palma} o la ontología de mantenimiento~\cite{woods}--- como de las ontologías previas del ciclo de vida de aprendizaje automático. Con las primeras comparte reutilización disciplinada, implementación ejecutable y evaluación contra preguntas de competencia, pero aplica ese patrón a la articulación entre catálogo de IA y plano de datos. Con las segundas comparte el dominio de aplicación, pero añade la capa de gobernanza necesaria para el intercambio en espacios de datos. Dos rasgos de diseño son especialmente relevantes: la verificación de implicación sobre el cierre OWL-RL, que complementa al razonamiento de consistencia, y las invariantes SHACL-SPARQL entre clases, evaluadas tanto sobre el grafo positivo como sobre un banco negativo.

Las limitaciones del perfil se agrupan en tres líneas argumentales. La primera concierne al \emph{nivel de detalle semántico} de varias extensiones DAIMO. Un límite transversal es el dinamismo del dominio: nuevas arquitecturas, formatos de despliegue, métricas y prácticas de evaluación evolucionan más rápido que los ciclos de mantenimiento de una ontología, por lo que deben entrar mediante versionado y perfiles complementarios antes que como ampliaciones precipitadas del núcleo. DAIMO cubre el descubrimiento de modelos de IA como activos gobernados: las CQs D1--D4 localizan modelos por tarea, dominio, política, endpoint y umbral de evaluación. No pretende ser un catálogo de catálogos ni un motor de consulta federada, ni cubrir por sí mismo el descubrimiento de modelos de datos, esquemas o vocabularios sectoriales. El \code{IOContract} captura el formato de entrada/salida y el método de autenticación; este último está restringido por \code{sh:in} a un vocabulario cerrado. Las referencias opcionales \code{inputSchema}, \code{outputSchema} y perfiles basados en \code{dct:conformsTo} pueden apuntar a esquemas o estándares externos, pero la compatibilidad semántica fina entre los esquemas del modelo y del consumidor queda fuera del núcleo y debe abordarse mediante perfiles complementarios, como \code{daimo-discovery}, o shapes específicas del espacio de datos. Del mismo modo, DAIMO describe ejecuciones autorizadas y auditables, no implementa enforcement federado: el uso de conectores, nodos neutrales, contenedores o TEE pertenece a la plataforma y puede reflejarse mediante metadatos o perfiles, no como axiomas universales. El \code{SharedEvaluationContext} reifica la combinación tarea-dataset-versión-protocolo-semilla suficiente para comparaciones homogéneas sobre una métrica única, pero no cubre protocolos heterogéneos de \emph{benchmarking}, conjuntos compuestos de métricas ni equivalencia matemática de indicadores calculados por procedimientos distintos. Las métricas de sesgo, equidad, seguridad o interpretabilidad pueden representarse como medidas de evaluación cuando existen y están contextualizadas; DAIMO no define una teoría universal de esos conceptos. Los atributos protegidos, técnicas de mitigación y auditorías de sesgo pertenecen a perfiles de cumplimiento o de dominio. El modelado de roles conflata el tipo (\code{ModelProvider}) con la asignación ---la tupla agente-rol-contexto-alcance---, conflación que OntoClean recomienda separar. MLDCAT-AP adopta la misma pragmática en \code{hasRegisteredUser}, lo que sitúa la contextualización fina en el horizonte del refinamiento futuro, no en el del error de diseño.

La segunda línea concierne a los \emph{compromisos de integración} y a los efectos que tienen sobre consumidores que combinen DAIMO con vocabularios adyacentes. La propiedad \code{daimo:records} declara \code{rdfs:range~prov:Activity}: bajo razonamiento RDFS, las evaluaciones MLDCAT-AP agrupadas dentro de un \code{CrossParticipantProvenanceRecord} son inferidas como actividades PROV aunque MLDCAT-AP no declare esa subsunción. El efecto es inocuo para consultas de auditoría, pero un consumidor que integre DAIMO, MLDCAT-AP y PROV-O puede obtener tipados sorprendentes sobre \code{it6:Evaluation}. En la misma línea, DAIMO no se ancla a una ontología de nivel superior (BFO, GFO, DOLCE o UFO): la ausencia responde a que un perfil de integración no es una ontología de referencia. Un trabajo futuro coherente sería acomodar la taxonomía a una de esas ontologías superiores sin renunciar a la reutilización directa de los vocabularios de dominio que motivan DAIMO.

Una variante práctica de este problema aparece en la soberanía de datos: puede ser necesario propagar restricciones desde un dataset de entrenamiento hacia el modelo o hacia sus artefactos derivados. DAIMO no axiomatiza esa propagación como regla universal, porque un modelo no hereda automáticamente todas las restricciones jurídicas o territoriales de sus datos de origen; la respuesta depende de licencia, contrato, jurisdicción, riesgo de memorización, técnica de entrenamiento y política acordada. Lo que DAIMO sí aporta es la estructura para representarla y auditarla cuando una organización la exige: \code{it6:trainedOn} y \code{daimo:contextDataset} conectan modelo, evaluación y dataset; ODRL expresa restricciones de uso; \code{ExecutionAuthorization} vincula esas condiciones con ejecuciones concretas; y la evidencia de auditoría junto con el registro de procedencia cruzada permiten reconstruir si el intercambio respetó la política. Reglas de propagación más estrictas deben implementarse como perfiles SHACL u ODRL específicos del espacio de datos, no como axiomas generales del núcleo.

La tercera línea concierne a las \emph{condiciones de reproducibilidad y validación externa}. Los resultados numéricos reportados dependen de un perfil de razonamiento concreto: el modo \code{inference=rdfs} en pyshacl, necesario para que las subpropiedades implicadas se materialicen antes de evaluar los \emph{shapes} de conformidad, y el cierre OWL-RL materializado antes de las consultas de competencia. Esta dependencia está documentada pero es una condición que un consumidor debe conocer para reproducir los resultados exactos. La evidencia actual es validación ejecutable del artefacto; no sustituye una revisión externa por expertos del dominio, de espacios de datos o de ingeniería ontológica. Una validación externa adecuada debería juzgar suficiencia, claridad, alineación con prácticas de espacios de datos y utilidad para MLOps o cumplimiento, manteniendo separadas esas observaciones de las CQs verificadas por SPARQL.

Conviene delimitar la relación entre DAIMO y MLDCAT-AP para disipar toda preocupación de duplicidad. El modelo propiamente dicho se reutiliza directamente mediante \code{it6:MachineLearningModel}, sin redefinición. La aportación de DAIMO se concentra en la articulación entre catálogo y plano de datos: las cuatro preguntas de competencia de la categoría de gobernanza ejercitan exclusivamente clases DAIMO-nativas y no pueden responderse sólo con MLDCAT-AP, pues requieren \code{AIAssetOffering} (CQ-G1), \code{ModelDeployment} (CQ-G2), \code{ExecutionAuthorization} (CQ-G3) y \code{CrossParticipantProvenanceRecord} (CQ-G4). Los dos perfiles son complementarios.

Esta delimitación también sitúa la relación con aplicaciones de descubrimiento como \code{ai-model-discovery}. DAIMO no debe absorber la lógica de ranking, recomendación, ejecución federada o interfaz de usuario de una aplicación concreta; esas capacidades pertenecen a una capa operacional que consume el grafo DAIMO, no al núcleo ontológico. Una evolución razonable sería definir un perfil complementario \code{daimo-discovery} que añada metadatos de indexación, señales agregadas de popularidad, latencia observada, huella computacional o sostenibilidad, siempre como extensiones derivables o asociadas a contextos de evaluación, sin modificar las catorce clases nativas del núcleo. De este modo, el proyecto \code{ai-model-discovery} puede actuar como aplicación futura y banco de explotación del perfil, mientras que DAIMO conserva un alcance publicable como ontología de integración para intercambio gobernado.

La misma regla aplica a las preguntas complementarias sobre disponibilidad, reproducibilidad profunda y confianza operacional. Varias de ellas ya están cubiertas por el perfil actual porque forman parte de la cadena catálogo-política-despliegue-ejecución; otras pertenecen al plano del activo MLDCAT-AP o a perfiles futuros. La Tabla~\ref{tab:future-profiles} resume la decisión de alcance. Esta tabla no introduce módulos ontológicos nuevos: documenta qué preguntas son núcleo DAIMO y cuáles deben evolucionar como perfiles especializados para no complejizar la contribución principal.

\begin{table*}[t]
\small\sf\centering
\caption{Cobertura de familias complementarias de preguntas y ruta de evolución. El criterio es mantener el núcleo DAIMO limitado al puente catálogo-dataspace y desplazar detalles de MLOps, seguridad o cumplimiento fino a perfiles especializados.\label{tab:future-profiles}}
\begin{tabular}{@{}p{3.2cm}p{5.9cm}p{6.0cm}@{}}
\toprule
\textbf{Familia} & \textbf{Cobertura en DAIMO actual} & \textbf{Ruta de extensión controlada} \\
\midrule
Disponibilidad del recurso & Endpoint mediante \code{ModelDeployment}, \code{dcat:DataService} y \code{dcat:endpointURL}; política y licencia mediante ODRL y \code{dct:license}; distribuciones descargables mediante DCAT cuando el modelo se publica como archivo; proveedor mediante roles DAIMO y contexto de participante. \code{dct:publisher} queda reservado para el mantenedor del registro o catálogo cuando proceda. & Mantener en el núcleo sólo la disponibilidad necesaria para descubrimiento e invocación. Formatos de artefacto como ONNX o PyTorch se expresan con \code{dct:format}/\code{dcat:Distribution} o con \code{IOContract} si afectan a la invocación. \\
Reproducibilidad técnica & Dataset de entrenamiento y evaluación mediante MLDCAT-AP, DCAT y PROV-O; artefactos de reproducibilidad mediante \code{it6:Flow}, evidencia, checksum y contexto compartido de evaluación. & Hiperparámetros, dependencias software, versiones de framework, GPU y memoria pertenecen a \code{daimo-reproducibility} o a reutilización más detallada de ML-Schema/MEX/MLDCAT-AP, no a nuevas clases núcleo. \\
Versionado y linaje de modelo & DAIMO puede representar derivación y revisión mediante PROV-O cuando el grafo lo declara, y puede auditar la cadena ejecución-artefacto-autorización. & Un perfil de linaje puede añadir reglas SHACL para \code{prov:wasDerivedFrom}, \code{prov:wasRevisionOf}, estado de obsolescencia y trazas de reentrenamiento sin alterar las catorce clases nativas. \\
Automatización MLOps & Fuera del núcleo. DAIMO define el contrato semántico que una herramienta puede poblar, pero no implementa conectores ni extracción automática. & Integradores desde Hugging Face, MLflow o registros de modelos pueden poblar grafos DAIMO y servir como validación aplicada posterior. \\
Monitorización y deriva & Fuera del núcleo salvo que los resultados se publiquen como evaluaciones contextualizadas. & Un perfil \code{daimo-monitoring} puede modelar deriva, degradación, latencia observada, huella computacional o sostenibilidad como métricas temporales o eventos de evaluación. \\
Seguridad y cumplimiento & El núcleo cubre política, autorización, evidencia de auditoría, checksum, firmante y timestamp. No modela vulnerabilidades, ataques o interpretación jurídica automática. & Perfiles \code{daimo-security} y \code{daimo-compliance} pueden añadir vulnerabilidades, firmas, controles AI Act, sesgo/equidad/explicabilidad y reglas SHACL/ODRL específicas de dominio. \\
\bottomrule
\end{tabular}
\end{table*}

Más allá del artefacto presente, DAIMO ilustra un patrón que probablemente se repetirá allí donde participantes autónomos necesiten intercambiar recursos de alto valor sujetos a política. El patrón tiene tres elementos recurrentes: un registro de oferta reificado sobre el propio recurso, una autorización reificada sobre el permiso, y un bundle de procedencia que cruza contextos administrativos. Conjuntos de datos genómicos, herramientas de apoyo a la decisión clínica, gemelos digitales y modelos industriales presentan cadenas análogas de publicación-invocación-auditoría. DAIMO es específica de IA en lo que reutiliza (MLDCAT-AP), pero agnóstica al dominio en los constructos que añade; confiamos en que esos constructos sirvan como punto de partida para perfiles hermanos en dominios adyacentes.

La ruta de adopción es concreta aunque el artículo no reporte todavía un piloto desplegado. En el plano del activo de IA, DAIMO puede coevolucionar con el dossier SEMIC/MLDCAT-AP preservando su capa de descripción de modelos y añadiendo únicamente el puente de gobernanza hacia espacios de datos. En el plano de intercambio, Eclipse Dataspace Components es el candidato natural de integración: los contextos de participante anclan procedencia agregada, y los acuerdos ODRL ya aparecen al final de una negociación DSP. En el plano de cumplimiento, la evidencia de auditoría DAIMO contiene los elementos que inspecciones alineadas con Gaia-X o con el Anexo~IV del AI Act esperarían encontrar: checksum, firmante, timestamp, autorización y vínculo con un bundle de procedencia cruzada. Estas observaciones definen tres rutas de validación posteriores: instanciación sobre un catálogo MLDCAT-AP público, co-despliegue con un conector EDC en un piloto federado y mapeo de un caso de alto riesgo del AI Act contra las seis invariantes SHACL-SPARQL.

% =====================================================================
\section{Disponibilidad y sostenibilidad}\label{sec:avail}

DAIMO se acompaña de los artefactos necesarios para inspeccionar y reproducir la evaluación: módulos OWL en Turtle, alineamientos, \emph{shapes} SHACL, grafo de demostración, consultas de competencia, banco negativo y reportes generados. La ontología y su documentación se declaran bajo licencia CC-BY~4.0, mientras que el código de validación se declara bajo Apache~2.0. La reproducibilidad no depende de un endpoint productivo: los resultados reportados pueden comprobarse localmente con las dependencias RDF, SHACL y OWL-RL indicadas por el paquete.

La sostenibilidad se plantea siguiendo la Fase~4 de LOT. Los módulos del perfil tienen IRIs versionadas, metadatos Dublin Core, política de versionado SemVer y una separación explícita entre núcleo, alineamientos y capa SHACL. Esta separación permite corregir documentación, actualizar alineamientos o extender perfiles especializados sin convertir cada evolución operativa en un cambio del núcleo. La publicación está orientada a los principios FAIR mediante identificadores persistentes, serializaciones RDF estándar, metadatos de licencia y definiciones SKOS por término; esos elementos se presentan como condiciones de disponibilidad y reutilización del artefacto, no como evidencia adicional de despliegue o adopción.

% =====================================================================
\section{Conclusiones}\label{sec:conclusion}

DAIMO conecta MLDCAT-AP, DCAT, ODRL, PROV-O y DSP para representar modelos de IA como activos gobernados dentro de un espacio de datos. Su contribución específica es la capa de gobernanza entre el catálogo de IA y el plano de datos: registro con política ODRL, autorización de ejecución anclada a un acuerdo, despliegue como puente entre modelo y servicio, y procedencia agregada entre contextos administrativos. El demostrador multi-participante muestra que el perfil puede instanciarse de forma coherente sobre un escenario controlado y responder las cinco familias de consultas operativas. Las líneas de trabajo futuro se derivan de las limitaciones declaradas: instanciación contra catálogos MLDCAT-AP reales, revisión cualitativa con expertos, refinamiento del \code{IOContract}, separación rol-tipo/rol-asignación, extensión del \code{SharedEvaluationContext} a protocolos heterogéneos y perfiles complementarios de reproducibilidad, monitorización, seguridad y cumplimiento. El artefacto se acompaña de una capa SHACL independiente y de una evaluación reproducible, sin reclamar despliegue productivo ni validación externa ya realizada.

% =====================================================================
\begin{thebibliography}{45}

\bibitem{franklin05} Franklin, M., Halevy, A. and Maier, D. (2005) From Databases to Dataspaces: A New Abstraction for Information Management. \emph{SIGMOD Record}, 34(4):27--33.

\bibitem{dcat2} Albertoni, A. et al. (2020) \emph{Data Catalog Vocabulary (DCAT) --- Version 2}. W3C Recommendation.

\bibitem{provo} Lebo, P., Sahoo, S. and McGuinness, D. (eds.) (2013) \emph{PROV-O: The PROV Ontology}. W3C Recommendation.

\bibitem{odrl22} Iannella, R., Guth, S. and Steyskal, R. (eds.) (2018) \emph{ODRL Information Model 2.2}. W3C Recommendation.

\bibitem{shacl} Knublauch, H. and Kontokostas, D. (eds.) (2017) \emph{Shapes Constraint Language (SHACL)}. W3C Recommendation.

\bibitem{mitchell19} Mitchell, M. et al. (2019) Model Cards for Model Reporting. \emph{Proceedings of the Conference on Fairness, Accountability, and Transparency}.

\bibitem{arnold19} Arnold, M. et al. (2019) FactSheets: Increasing Trust in AI Services through Supplier's Declarations of Conformity. \emph{IBM Research Report}.

\bibitem{werbrouck24} Werbrouck, J. et al. (2024) ConSolid: A federated ecosystem for heterogeneous multi-stakeholder projects. \emph{Semantic Web}, 15:429--460.

\bibitem{kurteva} Kurteva, A. et al. (in press) RePlanIT Ontology for Digital Product Passports of ICT: Laptops and Data Servers. \emph{Semantic Web}.

\bibitem{palma} Palma, R. et al. (in press) GloSIS: The Global Soil Information System Web Ontology. \emph{Semantic Web}.

\bibitem{woods} Woods, C. et al. (in press) An ontology for maintenance activities and its application to data quality. \emph{Semantic Web}.

\bibitem{penteado23} Penteado, B.E., Maldonado, J.C. and Isotani, S. (2023) Methodologies for publishing linked open government data on the Web. \emph{Semantic Web}, 14:585--610.

\bibitem{bareedu24} Bareedu, Y.S. et al. (2024) Deriving semantic validation rules from industrial standards: An OPC UA study. \emph{Semantic Web}, 15:517--554.

\bibitem{ferranti} Ferranti, N. et al. (in press) Formalizing and Validating Wikidata's Property Constraints using SHACL and SPARQL. \emph{Semantic Web}.

\bibitem{pandit-esteves} Pandit, H.J. and Esteves, B. (in press) Enhancing Data Use Ontology (DUO) for Health-Data Sharing by Extending it with ODRL and DPV. \emph{Semantic Web}.

\bibitem{robaldo} Robaldo, L. and Batsakis, S. (in press) On the interplay between validation and inference in SHACL. \emph{Semantic Web}.

\bibitem{otto19} Otto, B. et al. (2019) Reference architecture model for international data spaces. \emph{International Journal of Information Management}, 45:182--199.

\bibitem{cappiello22} Cappiello, C., Gal, A., Jarke, M., Rehof, J. and Yang, Y. (2022) Data spaces: Design, deployment, and future directions. \emph{ACM Computing Surveys}, 55(2):Article~46.

\bibitem{gebru21} Gebru, T. et al. (2021) Datasheets for datasets. \emph{Communications of the ACM}, 64(12):86--92.

\bibitem{kipoi} Avsec, Z. et al. (2019) The Kipoi repository accelerates community exchange and reuse of predictive models for genomics. \emph{Nature Biotechnology}, 37.

\bibitem{tfx} Baylor, D. et al. (2017) TFX: A TensorFlow-based production-scale machine learning platform. \emph{KDD}.

\bibitem{expo} Soldatova, L.N. and King, R.D. (2006) An ontology of scientific experiments. \emph{Journal of the Royal Society Interface}, 3(11).

\bibitem{mls} Correa Publio, G. et al. (2018) ML-Schema: Exposing the semantics of machine learning with schemas and ontologies.

\bibitem{mex} Keet, C.M., Fernández, J.D. and Panov, P. (2016) MEX vocabulary: A lightweight interchange format for machine learning experiments. \emph{ISWC}.

\bibitem{ontodm} Panov, P., Dzeroski, S. and Soldatova, L.N. (2008) OntoDM: An ontology of data mining. \emph{ICDMW}.

\bibitem{dmop} Keet, C.M. et al. (2015) The data mining optimization ontology. \emph{Web Semantics}, 32.

\bibitem{aio} Joachimiak, M.P. et al. (2024) The Artificial Intelligence Ontology: LLM-assisted construction of AI concept hierarchies. \emph{Applied Ontology}.

\bibitem{ito} Blagec, K., Barbosa-Silva, A., Ott, S. and Samwald, M. (2022) A curated, ontology-based, large-scale knowledge graph of artificial intelligence tasks and benchmarks. \emph{Scientific Data}, 9:322.

\bibitem{novacek24} Novacek, J., Ahari, A., M\"uller, T., Reiter, S., Viehl, A. and Bringmann, O. (2024) Ontology-Supported AI Model and Dataset Management. \emph{IEEE 22nd International Conference on Industrial Informatics (INDIN)}, 1--6.

\bibitem{souza19} Souza, R. et al. (2019) Provenance data in the machine learning lifecycle in computational science and engineering. \emph{WORKS}.

\bibitem{schmidt21} Schmidt, M. et al. (2021) FAIRness along the machine learning lifecycle using Dataverse in combination with MLflow. \emph{ICDE Workshops}.

\bibitem{lot22} Poveda-Villalón, M., Fernández-Izquierdo, A., Fernández-López, M. and García-Castro, R. (2022) LOT: An industrial oriented ontology engineering framework. \emph{Engineering Applications of Artificial Intelligence}, 111:104755.

\bibitem{brank05} Brank, J., Grobelnik, M. and Mladeni\'c, D. (2005) A survey of ontology evaluation techniques. \emph{Conference on Data Mining and Data Warehouses}.

\bibitem{oops14} Poveda-Villalón, M., Gómez-Pérez, A. and Suárez-Figueroa, M.C. (2014) OOPS! (OntOlogy Pitfall Scanner!). \emph{International Journal on Semantic Web and Information Systems}, 10(2):7--34.

\bibitem{mldcatap} SEMIC (2025) \emph{MLDCAT-AP 3.0.0: Machine Learning Data Catalogue Application Profile}. European Commission.

\bibitem{mldcatap31} SEMIC (2026) \emph{MLDCAT-AP 3.1.0: Machine Learning Data Catalogue Application Profile}. European Commission, Candidate Recommendation, 13 May 2026.

\bibitem{dsp} International Data Spaces Association (2024) \emph{Dataspace Protocol (DSP)}. W3C Community Draft.

\bibitem{dsp2025} Eclipse Dataspace Protocol (2025) \emph{Dataspace Protocol Release 2025-1}. Eclipse Dataspace Working Group.

\bibitem{edc} Eclipse Foundation. \emph{Eclipse Dataspace Components (EDC)}, control-plane documentation.

\bibitem{widoco} Garijo, D. (2017) WIDOCO: A Wizard for Documenting Ontologies. \emph{ISWC}.

\bibitem{aiact} European Union (2024) \emph{Regulation (EU) 2024/1689 laying down harmonised rules on artificial intelligence (AI Act)}. Official Journal of the European Union.

\end{thebibliography}

% =====================================================================
\appendix

\section{Preguntas de competencia}\label{app:cqs}

\subsection*{R --- Registro y publicación (5)}
\begin{enumerate}[label=\arabic*., leftmargin=*]
\item \cq{CQ-R1}: ¿Qué modelos ha publicado un proveedor dado como activos de catálogo gobernados?
\item \cq{CQ-R2}: Dado un modelo, ¿qué oferta lo registró, por qué proveedor y en qué fecha? (requiere cadena \code{subPropertyOf})
\item \cq{CQ-R3}: ¿Qué licencia y política de uso aplican a un modelo publicado?
\item \cq{CQ-R4}: ¿Qué contrato de entrada/salida se declara para la interfaz de invocación?
\item \cq{CQ-R5}: Para cada oferta, ¿qué subclase concreta de \code{ParticipantRole} tiene el agente proveedor? (requiere \code{subClassOf+})
\end{enumerate}

\subsection*{D --- Descubrimiento y selección (4)}
\begin{enumerate}[label=\arabic*., leftmargin=*]
\item \cq{CQ-D1}: ¿Qué modelos resuelven una tarea dada en un dominio dado?
\item \cq{CQ-D2}: ¿Qué modelos son utilizables bajo un patrón de política concreto (por ejemplo uso no comercial)?
\item \cq{CQ-D3}: ¿Qué modelos exponen un endpoint y qué método de autenticación requiere cada uno?
\item \cq{CQ-D4}: ¿Qué modelos alcanzan un umbral mínimo de métrica bajo un contexto de evaluación dado?
\end{enumerate}

\subsection*{E --- Ejecución y auditabilidad (5)}
\begin{enumerate}[label=\arabic*., leftmargin=*]
\item \cq{CQ-E1}: Dada una oferta, ¿qué endpoint, método de autenticación y contrato de entrada/salida aplican al modelo subyacente? (requiere entailment)
\item \cq{CQ-E2}: Para un despliegue dado, ¿qué ejecuciones, por qué agentes y bajo qué autorización?
\item \cq{CQ-E3}: Para una ejecución dada, ¿qué implementación, algoritmo, \emph{flow} e infraestructura?
\item \cq{CQ-E4}: ¿Qué evidencia de auditoría (hash, firmante, marca temporal) respalda una ejecución?
\item \cq{CQ-E5}: ¿Qué artefactos derivados produjo una ejecución y bajo qué autorización?
\end{enumerate}

\subsection*{V --- Evaluación y reproducibilidad (5)}
\begin{enumerate}[label=\arabic*., leftmargin=*]
\item \cq{CQ-V1}: ¿Qué contexto compartido (dataset, versión, protocolo, semilla) usa una evaluación?
\item \cq{CQ-V2}: Bajo un contexto compartido, ¿qué modelo alcanza el mayor valor de métrica?
\item \cq{CQ-V3}: ¿Cómo se ordenan dos o más modelos bajo el mismo contexto y métrica?
\item \cq{CQ-V4}: Para un modelo dado, ¿sobre qué benchmark suites ha sido evaluado?
\item \cq{CQ-V5}: ¿Qué artefactos de reproducibilidad (flow, cuaderno, tabla de resultados, checksum) respaldan una evaluación?
\end{enumerate}

\subsection*{G --- Puente de gobernanza (4, DAIMO-nativas)}
\begin{enumerate}[label=\arabic*., leftmargin=*]
\item \cq{CQ-G1}: ¿Qué ofertas del catálogo federado incluyen un modelo dado?
\item \cq{CQ-G2}: ¿Qué despliegues sirven un modelo dado, sobre qué infraestructura y con qué contrato de entrada/salida?
\item \cq{CQ-G3}: ¿Qué autorización, y el acuerdo del que deriva, permitió una ejecución específica?
\item \cq{CQ-G4}: A través de contextos de participante, ¿cuál es el bundle de procedencia completo para un artefacto derivado?
\end{enumerate}

\section{Shapes SHACL}\label{app:shapes}

La capa SHACL de DAIMO contiene nueve \emph{shapes} nodales de completitud, tres \emph{shapes} de conformidad sobre clases reutilizadas y seis invariantes SHACL-SPARQL entre clases. El módulo \code{shapes/daimo-shapes.ttl} se declara a su vez como \code{owl:Ontology} independiente con \code{owl:versionIRI} bajo \url{https://w3id.org/pionera/daimo/shapes/}.

\begin{itemize}[leftmargin=*]
\item \textbf{Shapes de completitud (9)}: \code{AIAssetOfferingShape}, \code{ParticipantRoleShape}, \code{ModelDeploymentShape}, \code{IOContractShape}, \code{ExecutionAuthorizationShape}, \code{DerivedArtifactShape}, \code{SharedEvaluationContextShape}, \code{AuditEvidenceShape} y \code{CrossParticipantProvenanceRecordShape}.
\item \textbf{Shapes de conformidad sobre clases reutilizadas (3)}: \code{OfferInDAIMOShape} exige \code{odrl:assigner} y \code{odrl:target} sobre toda \code{odrl:Offer}; \code{MachineLearningModelInDAIMOShape} exige política, título e identificador sobre \code{it6:MachineLearningModel}; \code{RunInDAIMOShape} exige \emph{flow}, algoritmo, agente asociado y marca temporal sobre \code{it6:Run}.
\item \textbf{Invariantes cruzadas (6)}: \code{DerivationAuthorizationInvariant} (INV-1), \code{RunAgentAuthorizationInvariant} (INV-2), \code{DeploymentServiceInvariant} (INV-3), \code{AuthorizationTemporalInvariant} (INV-4), \code{OfferingPolicyTargetInvariant} (INV-5) y \code{OfferingAssignerInvariant} (INV-6); todas como \code{sh:SPARQLConstraint}.
\end{itemize}

\section{Consultas SPARQL}\label{app:sparql}

Las veintitrés consultas SPARQL correspondientes a las preguntas del Apéndice~\ref{app:cqs} están publicadas en \code{daimo/queries/queries.md}. El validador ejecuta cada consulta sobre el cierre OWL-RL del grafo unión ontología más ejemplo y verifica que cada una devuelve al menos una fila. Los conteos exactos aparecen en la Tabla~\ref{tab:cqs}. CQ-G2 devuelve dos filas ---una por cada endpoint del despliegue multi-protocolo REST y gRPC---, lo que confirma que el modelado de múltiples servicios por despliegue es un patrón legítimo y directamente observable desde SPARQL.

\section{Glosario de abreviaturas}\label{app:glossary}

\begin{table*}[h]
\small\sf\centering
\caption{Abreviaturas utilizadas en el artículo.\label{tab:glossary}}
\begin{tabular}{@{}p{2.6cm}p{12.2cm}@{}}
\toprule
\textbf{Abreviatura} & \textbf{Significado} \\
\midrule
AI Act & Regulation (EU) 2024/1689 laying down harmonised rules on artificial intelligence. \\
CQ & Competency Question. \\
DAIMO & Data-space AI Model Ontology. \\
DCAT & Data Catalog Vocabulary. \\
DCAT-AP & DCAT Application Profile. \\
DL & Description Logic. \\
DOI & Digital Object Identifier. \\
DPV & Data Privacy Vocabulary. \\
DSP & Dataspace Protocol. \\
EDC & Eclipse Dataspace Components. \\
FAIR & Findable, Accessible, Interoperable, Reusable. \\
LOT & Linked Open Terms. \\
MLDCAT-AP & Machine Learning Data Catalogue Application Profile. \\
ODRL & Open Digital Rights Language. \\
OOPS! & OntOlogy Pitfall Scanner. \\
OWL & Web Ontology Language. \\
PROV-O & PROV Ontology. \\
RDF & Resource Description Framework. \\
RDFS & RDF Schema. \\
SHACL & Shapes Constraint Language. \\
SKOS & Simple Knowledge Organization System. \\
SPARQL & SPARQL Protocol and RDF Query Language. \\
SWJ & Semantic Web Journal. \\
TTL & Turtle. \\
WIDOCO & Wizard for Documenting Ontologies. \\
\bottomrule
\end{tabular}
\end{table*}

\section{Convenciones de URI}\label{app:uris}

DAIMO utiliza el namespace canónico \url{https://w3id.org/pionera/daimo}, con identificadores de términos en la forma \code{https://w3id.org/pionera/daimo\#<Término>} e IRIs versionadas bajo \code{https://w3id.org/pionera/daimo/<versión>}, enlazadas entre sí mediante \code{owl:priorVersion}. Los módulos satélite son ontologías independientes con IRIs versionadas propias: la capa SHACL bajo \url{https://w3id.org/pionera/daimo/shapes/} y el módulo de alineamientos bajo \url{https://w3id.org/pionera/daimo/align} (este último declarando \code{owl:imports} sobre la ontología base). El grafo de ejemplo utiliza el namespace local \code{https://example.org/daimo-scenario/}, claramente separado del espacio canónico. La negociación de contenido resuelve el mismo identificador a HTML, Turtle, RDF/XML, JSON-LD o N-Triples en función del \code{Accept} header; la configuración \code{.htaccess} correspondiente está preparada en \code{w3id-redirect/.htaccess} a la espera de su integración en \code{perma-id/w3id.org}.

La elección de IRIs con \code{\#} se limita a términos de vocabulario, no a instancias operativas de alto volumen. En un espacio de datos productivo, modelos concretos, datasets, ejecuciones y artefactos derivados deben identificarse bajo namespaces de los participantes o del catálogo federado; DAIMO sólo proporciona los predicados y clases con los que esos recursos se describen. Los términos nativos usan inglés ASCII: clases en \emph{UpperCamelCase} singular (\code{AIAssetOffering}, \code{ModelDeployment}), propiedades en \emph{lowerCamelCase} verbal o relacional (\code{offersModel}, \code{authorizesRun}, \code{usesEvaluationContext}) y módulos/versiones mediante segmentos \emph{slash}. Las traducciones, sinónimos y explicaciones se expresan con \code{rdfs:label}, \code{rdfs:comment} y SKOS, no mediante variantes de URI.

\end{document}
